% Options for packages loaded elsewhere
\PassOptionsToPackage{unicode}{hyperref}
\PassOptionsToPackage{hyphens}{url}
\PassOptionsToPackage{dvipsnames,svgnames,x11names}{xcolor}
\documentclass[
  american,
  a4paper,
]{article}
\usepackage{xcolor}
\usepackage[margin=25mm]{geometry}
\usepackage{amsmath,amssymb}
\setcounter{secnumdepth}{-\maxdimen} % remove section numbering
\usepackage{iftex}
\ifPDFTeX
  \usepackage[T1]{fontenc}
  \usepackage[utf8]{inputenc}
  \usepackage{textcomp} % provide euro and other symbols
\else % if luatex or xetex
  \usepackage{unicode-math} % this also loads fontspec
  \defaultfontfeatures{Scale=MatchLowercase}
  \defaultfontfeatures[\rmfamily]{Ligatures=TeX,Scale=1}
\fi
\usepackage{lmodern}
\ifPDFTeX\else
  % xetex/luatex font selection
  \ifLuaTeX
    \usepackage{luaotfload}
    \directlua{luaotfload.add_fallback("mainfontfallback",{
      "DejaVu Serif:","DejaVu Sans:"
    })}
  \fi
  \setmainfont[,RawFeature={fallback=mainfontfallback}]{Latin Modern Roman}
  \setmonofont[]{DejaVu Sans Mono}
\fi
% Use upquote if available, for straight quotes in verbatim environments
\IfFileExists{upquote.sty}{\usepackage{upquote}}{}
\IfFileExists{microtype.sty}{% use microtype if available
  \usepackage[]{microtype}
  \UseMicrotypeSet[protrusion]{basicmath} % disable protrusion for tt fonts
}{}
\makeatletter
\@ifundefined{KOMAClassName}{% if non-KOMA class
  \IfFileExists{parskip.sty}{%
    \usepackage{parskip}
  }{% else
    \setlength{\parindent}{0pt}
    \setlength{\parskip}{6pt plus 2pt minus 1pt}}
}{% if KOMA class
  \KOMAoptions{parskip=half}}
\makeatother
\usepackage{longtable,booktabs,array}
\usepackage{calc} % for calculating minipage widths
% Correct order of tables after \paragraph or \subparagraph
\usepackage{etoolbox}
\makeatletter
\patchcmd\longtable{\par}{\if@noskipsec\mbox{}\fi\par}{}{}
\makeatother
% Allow footnotes in longtable head/foot
\IfFileExists{footnotehyper.sty}{\usepackage{footnotehyper}}{\usepackage{footnote}}
\makesavenoteenv{longtable}
\usepackage{graphicx}
\makeatletter
\newsavebox\pandoc@box
\newcommand*\pandocbounded[1]{% scales image to fit in text height/width
  \sbox\pandoc@box{#1}%
  \Gscale@div\@tempa{\textheight}{\dimexpr\ht\pandoc@box+\dp\pandoc@box\relax}%
  \Gscale@div\@tempb{\linewidth}{\wd\pandoc@box}%
  \ifdim\@tempb\p@<\@tempa\p@\let\@tempa\@tempb\fi% select the smaller of both
  \ifdim\@tempa\p@<\p@\scalebox{\@tempa}{\usebox\pandoc@box}%
  \else\usebox{\pandoc@box}%
  \fi%
}
% Set default figure placement to htbp
\def\fps@figure{htbp}
\makeatother
\usepackage{svg}
% definitions for citeproc citations
\NewDocumentCommand\citeproctext{}{}
\NewDocumentCommand\citeproc{mm}{%
  \begingroup\def\citeproctext{#2}\cite{#1}\endgroup}
\makeatletter
 % allow citations to break across lines
 \let\@cite@ofmt\@firstofone
 % avoid brackets around text for \cite:
 \def\@biblabel#1{}
 \def\@cite#1#2{{#1\if@tempswa , #2\fi}}
\makeatother
\newlength{\cslhangindent}
\setlength{\cslhangindent}{1.5em}
\newlength{\csllabelwidth}
\setlength{\csllabelwidth}{3em}
\newenvironment{CSLReferences}[2] % #1 hanging-indent, #2 entry-spacing
 {\begin{list}{}{%
  \setlength{\itemindent}{0pt}
  \setlength{\leftmargin}{0pt}
  \setlength{\parsep}{0pt}
  % turn on hanging indent if param 1 is 1
  \ifodd #1
   \setlength{\leftmargin}{\cslhangindent}
   \setlength{\itemindent}{-1\cslhangindent}
  \fi
  % set entry spacing
  \setlength{\itemsep}{#2\baselineskip}}}
 {\end{list}}
\usepackage{calc}
\newcommand{\CSLBlock}[1]{\hfill\break\parbox[t]{\linewidth}{\strut\ignorespaces#1\strut}}
\newcommand{\CSLLeftMargin}[1]{\parbox[t]{\csllabelwidth}{\strut#1\strut}}
\newcommand{\CSLRightInline}[1]{\parbox[t]{\linewidth - \csllabelwidth}{\strut#1\strut}}
\newcommand{\CSLIndent}[1]{\hspace{\cslhangindent}#1}
\ifLuaTeX
\usepackage[bidi=basic,shorthands=off,]{babel}
\else
\usepackage[bidi=default,shorthands=off,]{babel}
\fi
\ifPDFTeX
\else
\babelfont{rm}[,RawFeature={fallback=mainfontfallback}]{Latin Modern Roman}
\fi
\ifLuaTeX
  \usepackage{selnolig} % disable illegal ligatures
\fi
\setlength{\emergencystretch}{3em} % prevent overfull lines
\providecommand{\tightlist}{%
  \setlength{\itemsep}{0pt}\setlength{\parskip}{0pt}}
\let\oldsection\section\renewcommand{\section}{\clearpage\oldsection}
\usepackage{bookmark}
\IfFileExists{xurl.sty}{\usepackage{xurl}}{} % add URL line breaks if available
\urlstyle{same}
\hypersetup{
  pdflang={en-US},
  colorlinks=true,
  linkcolor={black},
  filecolor={Maroon},
  citecolor={black},
  urlcolor={blue},
  pdfcreator={LaTeX via pandoc}}

\author{}
\date{}

\begin{document}

% HSE DSBA formal title page — matches final_report_coursework_3.pdf format.
% Included by pandoc via --include-before-body=reports/titlepage.tex.
% Cyrillic group codes (БПАД, БПМИ) render via the DejaVu Serif fallback
% configured in the pandoc invocation (see KB/build.md).
\begin{titlepage}
\thispagestyle{empty}
\centering

{\Large\bfseries NATIONAL RESEARCH UNIVERSITY\par}
\vspace{0.25cm}
{\Large\bfseries HIGHER SCHOOL OF ECONOMICS\par}

\vspace{1.8cm}

{\large Faculty of Computer Science\par}
\vspace{0.3cm}
{\large Bachelor's Programme ``Data Science and Business Analytics''\par}

\vspace{2.2cm}

{\large Software Project Report on the Topic:\par}
\vspace{0.6cm}
{\LARGE\bfseries Neural Networks Adaptation for Predictive\\
Diagnostics Domain in Medicine\par}

\vspace{2.2cm}

\begin{flushleft}
\hspace{1cm}\textbf{Fulfilled by:}\par
\vspace{0.3cm}
\hspace{1cm}\begin{minipage}{0.92\textwidth}
Student of the Group БПАД232\hfill Suvorova Aleksandra Viacheslavovna\par
\vspace{0.05cm}
\hfill\rule{4cm}{0.3pt}\par
\vspace{-0.15cm}
\hfill (date)\par
\vspace{0.5cm}
Student of the Group БПМИ228\hfill Murzova Sofia\par
\vspace{0.05cm}
\hfill\rule{4cm}{0.3pt}\par
\vspace{-0.15cm}
\hfill (date)\par
\end{minipage}

\vspace{0.7cm}

\hspace{1cm}\textbf{Assessed by the Project Supervisor:}\par
\vspace{0.3cm}
\hspace{1cm}\begin{minipage}{0.92\textwidth}
Borevskiy Andrey Olegovich\par
Assistant\par
Department of Data Analysis and AI, CS HSE\par
\end{minipage}
\end{flushleft}

\vfill

{\large Moscow 2026\par}

\vspace{0.5cm}

\end{titlepage}

\renewcommand*\contentsname{Contents}
{
\setcounter{tocdepth}{2}
\tableofcontents
}
\section{Abstract}\label{abstract}

Clinical decision support requires retrieval grounded in vetted
evidence, since general-purpose language models fabricate dosages,
diagnostic criteria, and clinical statistics when their parametric
memory is queried directly (\citeproc{ref-singhal2023medpalm}{Singhal et
al., 2023}). This prototype implements a multi-agent medical RAG system
covering four specialties --- cardiology, endocrinology,
gastroenterology, and infectious diseases --- with three architectural
components: an LLM router that classifies the clinical domain of each
query, per-specialty FAISS dense and BM25 sparse retrieval over a
2,561-document corpus, and a multi-judge faithfulness evaluator
combining two same-family judges with one cross-vendor judge.

Evaluation used a 200-case golden dataset stratified into core,
peripheral, and out-of-scope tiers, with a 60-case development split for
hyperparameter tuning and a 140-case held-out test split for the
reported numbers. On the test split the LLM router reached 95.7\%
accuracy across the four specialties; FAISS dense retrieval reached
Recall@5 (the fraction of gold-source documents retrieved in the top-5
window) of 56.2\% averaged over per-specialty gold-document Bernoulli
trials; minimum-judge faithfulness (a case counts FAITHFUL only when
every configured judge agrees) reached 96.2\% under the three-judge
cross-vendor panel and 99.2\% under the two same-family judges, with the
−3.0 pp gap tracking the same-family inflation signal Zheng et al.
(\citeproc{ref-zheng2023mtbench}{2023}) predict; and the numeric
out-of-scope refusal gate caught 20.7\% of Tier 3 queries at a 17.1\%
Tier 1/2 false-positive rate.

The refusal gate is the principal limitation: both target ratios (≥80\%
Tier 3 recall, ≤5\% Tier 1/2 FP) are unmet because the in-corpus and
out-of-corpus L2 distributions overlap heavily at four-specialty scale.
The 56.2\% retrieval headline also sits below the dense-only ceiling
reported by recent medical RAG benchmarks
(\citeproc{ref-xiong2024medrag}{Xiong et al., 2024}); a two-stage gate
and a hybrid retriever are the immediate follow-ups.

\subsection{Keywords}\label{keywords}

medical RAG; retrieval-augmented generation; multi-agent system;
LLM-as-judge; faithfulness evaluation; refusal gate; FAISS; BM25

\section{1. Introduction}\label{introduction}

Clinical decision support systems require accurate, domain-specific
evidence to function safely. While general-purpose Large Language Models
(LLMs) can process complex medical queries, their parametric memory is
prone to fabricating critical medical facts such as dosages, diagnostic
criteria, and clinical statistics --- a failure mode quantified for
clinical QA by Singhal et al. (\citeproc{ref-singhal2023medpalm}{2023})
and for medical RAG specifically by Xiong et al.
(\citeproc{ref-xiong2024medrag}{2024}). A multi-agent
Retrieval-Augmented Generation (RAG) architecture offers a candidate
solution by forcing the LLM to ground its reasoning exclusively in
verified medical literature retrieved from specialist-specific vector
indices. This report evaluates such a prototype, designed for academic
investigation rather than immediate clinical use.

This work contributes the following:

\begin{enumerate}
\def\labelenumi{\arabic{enumi}.}
\item
  A four-specialty multi-agent medical RAG architecture evaluated
  end-to-end on a 200-case tiered golden dataset, with per-specialty
  FAISS+BM25 retrieval and an LLM router that achieves 95.7\% test-split
  routing accuracy.
\item
  A multi-judge cross-vendor faithfulness protocol that bounds
  same-family-LLM-judge bias to a measured −3.0 pp constraint, with
  upper-bound (99.2\%) and lower-bound (96.2\%) faithfulness rates
  reported on the held-out test split.
\item
  A characterisation of the architectural insufficiency of
  single-threshold L2 gating for out-of-scope refusal at four-specialty
  scale --- the paper's principal finding. The threshold sweep shows no
  scalar \texttt{L2\_REJECT\_MIN} value simultaneously meets both target
  rates (≥80\% Tier 3 catch, ≤5\% Tier 1/2 false positive); the chosen
  threshold delivers 20.7\% catch with a measured 17.1\% false-positive
  rate, motivating a two-stage L2-plus-classifier gate as the next
  architectural step.
\end{enumerate}

\begin{enumerate}
\def\labelenumi{(\arabic{enumi})}
\tightlist
\item
  Does an LLM-based query router add measurable clinical value over a
  deterministic keyword-matching baseline? (2) How does vector retrieval
  quality degrade when moving from core textbook conditions to
  peripheral or out-of-scope clinical scenarios? (3) Can an LLM acting
  as a strict faithfulness judge reliably detect medical hallucinations
  in generated responses?
\end{enumerate}

All headline numbers below are from the held-out test split (n = 140).
The principal unresolved limitation is the numeric out-of-scope refusal
gate: at the re-tuned threshold \texttt{L2\_REJECT\_MIN\ =\ 1.020} it
catches only \textbf{20.7\% (6/29) of Tier 3 cases at a 17.1\% Tier 1/2
false-positive rate}, leaving the original ≥80\% recall / ≤5\% FP
targets unmet, and a two-stage gate (numeric pre-filter + LLM confirmer)
is the proposed remediation (§5.5, §7 L7).

Against that backdrop, the LLM router achieves \textbf{95.7\% accuracy
(134/140) {[}Wilson 95\% CI 91.0\%--98.0\%{]}} across four specialties
(cardiology, endocrinology, gastroenterology, infectiology); the 6
misses are all defensible cross-specialty cases (§5.1). FAISS dense
retrieval reaches \textbf{Recall@5 of 56.2\% (162/288 gold-doc trials)
{[}50.5\%--61.9\%{]}} --- 22.2 pp above the BM25 sparse baseline
(34.0\%) and 43.8 pp short of an oracle ceiling (§5.3, §5.3.2).
Minimum-judge faithfulness is bounded above by the 2-Yandex-judge rate
of \textbf{99.2\% (131/132) {[}95.8\%--99.9\%{]}} (single disagreement:
\texttt{cardio\_40}, Tier 2 cardiology --- congenital long QT syndrome;
unchanged from the prior 2-specialty baseline) and bounded below by the
3-judge cross-vendor rate of \textbf{96.2\% (127/132)
{[}91.4\%--98.4\%{]}}, where the OpenAI GPT-OSS-120B tertiary judge
added 4 HALLUCINATION verdicts on cases both Yandex judges marked
FAITHFUL --- the empirical −3.0 pp same-family-bias signal Zheng et al.
(\citeproc{ref-zheng2023mtbench}{2023}) predict (§5.4).

The principal direction beyond this prototype is the two-stage refusal
gate described in §7 L7: pairing the L2 distance signal with a
generation-step claim-coverage check is the most concrete route to
closing the gap between the 20.7\% Tier 3 catch and a target rate that
would justify clinical use.

\subsection{Objectives}\label{objectives}

\begin{enumerate}
\def\labelenumi{\arabic{enumi}.}
\tightlist
\item
  Route clinical queries to the correct specialist agent with high
  accuracy.
\item
  Retrieve contextually relevant evidence from a large medical corpus.
\item
  Generate responses that are faithful to the retrieved evidence, with
  no medical hallucinations (fabricated drug names, dosages, diagnostic
  criteria, or statistics).
\end{enumerate}

\begin{center}\rule{0.5\linewidth}{0.5pt}\end{center}

\section{2. Related Work}\label{related-work}

Four research strands directly shape this prototype's design: medical
RAG benchmarks, retrieval-method baselines, LLM-as-judge evaluation
methodology, and multi-agent clinical-AI architectures.

\subsection{2.1 Medical RAG benchmarks}\label{medical-rag-benchmarks}

Early medical QA benchmarks measured factual recall on closed-form
questions; recent work measures retrieval quality and generation
faithfulness jointly on specialist corpora. Singhal et al.'s Med-PaLM
(\citeproc{ref-singhal2023medpalm}{Singhal et al., 2023}) introduced
MultiMedQA, a composite of six clinical and consumer-health QA datasets
(MedQA, PubMedQA, MedMCQA, LiveQA, MedicationQA, HealthSearchQA), and
showed that domain-instruction-tuned LLMs can match human expert
preference on a substantial fraction of consumer questions while still
fabricating dangerous specifics in others. Xiong et al.'s MedRAG/MIRAGE
benchmark (\citeproc{ref-xiong2024medrag}{Xiong et al., 2024}) extends
this to RAG specifically: 7,663 multiple-choice questions across five
biomedical corpora, with explicit retrieval/recall/precision per
question and an analysis of how chunk granularity and retriever choice
affect downstream accuracy. BioASQ
(\citeproc{ref-tsatsaronis2015bioasq}{Tsatsaronis et al., 2015})
predates both and provides the largest continuously-curated biomedical
semantic-indexing + QA shared task, with expert-annotated relevance
judgements and ideal/exact-answer pairs over 10+ years of editions.

Unlike MIRAGE and BioASQ, this evaluation uses open-ended clinical
scenarios rather than multiple-choice items, stratifies cases by
difficulty tier, and explicitly tests refusal on out-of-scope queries.
The trade-off is reduced comparability with prior medical-RAG numbers;
the comparable surface is the retrieval-quality methodology (top-K, Hit
Rate, Recall@K with gold-source labels) and the LLM-as-judge
faithfulness protocol described below.

\subsection{2.2 Retrieval method
baselines}\label{retrieval-method-baselines}

The canonical RAG framework was introduced by Lewis et al.
(\citeproc{ref-lewis2020rag}{2020}): a dense retriever feeds top-K
passages to a generator, both jointly fine-tuned end-to-end. Dense
Passage Retrieval (\citeproc{ref-karpukhin2020dpr}{Karpukhin et al.,
2020}) is the dominant dense baseline --- a dual-encoder learned via
in-batch contrastive loss on Natural Questions / TriviaQA --- and the
canonical sparse baseline is BM25
(\citeproc{ref-robertson2009bm25}{Robertson \& Zaragoza, 2009}), a
probabilistic IDF-weighted lexical match.

This project uses \textbf{dense-only Yandex
\texttt{text-search-doc}/\texttt{text-search-query} asymmetric
embeddings}, with neither a BM25 sparse baseline nor a hybrid fusion
(§5.3.2 evaluates BM25 as a comparison ablation rather than a production
component). The decision was a deliberate scope reduction in an earlier
dense-only configuration: chunk-size tuning pushed the dev-set Hit Rate
to 96.7\%; keyword stripping contributed no additional improvement (see
§4.5). With this evaluation's grounded Recall@K (58.0\% full-set
Recall@K vs the legacy 87.5\% KeywordHitRate, §5.3 / §5.7) it is now
clear that the dense-only choice does miss a substantial fraction of
relevant documents. As §5.3.2 confirms, BM25 recovers rare entity names
that FAISS misses; hybrid score fusion is the natural next architecture,
though it is not evaluated here. The
\texttt{metadata{[}\textquotesingle{}keywords\textquotesingle{}{]}}
field on every chunk is preserved precisely for that BM25 future use.

\subsection{2.3 LLM-as-judge
methodology}\label{llm-as-judge-methodology}

RAGAS (\citeproc{ref-es2023ragas}{Es et al., 2023}) formalised automated
RAG evaluation with three LLM-judged metrics --- faithfulness, answer
relevance, and context precision --- by asking a strong LLM to score
each generated answer against the retrieved context. Zheng et al.'s
MT-Bench / Chatbot Arena work (\citeproc{ref-zheng2023mtbench}{Zheng et
al., 2023}) systematically characterised LLM-judge biases: position
bias, verbosity bias, and \textbf{same-family self-preference bias} ---
a judge LLM tends to rate outputs from its own model family more
favourably than outputs from other families on the same task. Their
measured magnitude (model-pair-dependent, but typically a 5--25
percentage-point inflation) directly motivates this project's two-judge
protocol: instead of trusting the single YandexGPT judge that produces a
100\% faithfulness rate on test, a second YandexGPT-Lite judge is
deployed with the same prompt and the \textbf{minimum-judge rate} is
reported (a case is FAITHFUL only when both judges agree).

On 2026-05-23 a third, \textbf{cross-vendor} judge --- OpenAI's
open-weights GPT-OSS-120B served via OpenRouter --- was added as a
cross-vendor tertiary judge; this third judge substantially bounds the
residual same-family bias. The 3-judge minimum-judge rate is −3.0 pp
below the 2-Yandex-judge rate, the empirical signal they predict (§5.4).

\subsection{2.4 Multi-agent medical
systems}\label{multi-agent-medical-systems}

Kim et al.'s MDAgents (\citeproc{ref-kim2024mdagents}{Kim et al., 2024})
is the closest recent multi-agent clinical-AI work. It builds a
\emph{collaboration} of LLMs that adaptively choose between solo,
paired, or group-discussion modes depending on the medical query's
complexity --- modelled on how human clinicians escalate from
single-physician to multi-disciplinary-team review. Each MDAgent's role
is dynamically assigned per case (radiologist, pathologist, clinician,
etc.) and the agents iteratively \emph{debate} the diagnosis, with the
framework choosing the level of collaboration based on internal
complexity estimates.

This project is \textbf{not} a multi-agent system in the MDAgents sense.
The ``multi-agent'' label here refers to a \textbf{single-step routing
architecture}: the orchestrator picks exactly one specialist agent per
query (cardiologist \emph{or} endocrinologist) and that agent answers in
isolation; there is no inter-agent communication, no debate, no role
re-assignment per case. Single-agent routing keeps retrieval sharp and
the routing decision cleanly evaluable, though the \texttt{cardio\_40}
case illustrates the ceiling: a presentation spanning cardiology and
medical genetics gets a cardiology-only answer where a
multi-disciplinary discussion would arguably serve the patient better.
Extending this prototype toward an MDAgents-style collaboration is
recorded as a future-work direction in §8; doing so would require
re-architecting both retrieval (cross-corpus search) and faithfulness
evaluation (multi-agent answer fusion).

\begin{center}\rule{0.5\linewidth}{0.5pt}\end{center}

\section{3. System Architecture}\label{system-architecture}

\begin{verbatim}
User Query
    |
    v
+--------------------------------------------------+
|  Orchestrator                                    |
|                                                  |
|  1. Safety Gate (regex)                          |
|     • Emergency phrases → immediate disclaimer   |
|     • Prescription requests → refusal            |
|                                                  |
|  2. LLM Router                                   |
|     • Model: YandexGPT/latest                    |
|     • temperature=0.0, max_tokens=10             |
|     • Output: one-word specialist name           |
+----------------------+---------------------------+
                       |
   +-----------+-------+--------+--------------+
   v           v                v              v
+--------+ +--------+      +---------+    +---------+
| Cardio | | Endo   |      | Gastro  |    | Infect  |
| Agent  | | Agent  |      | Agent   |    | Agent   |
+---+----+ +---+----+      +----+----+    +----+----+
    |         |                |              |
    v         v                v              v
 FAISS     FAISS            FAISS          FAISS
 7,730     37,791           8,670          7,476
 vectors   vectors          vectors        vectors
    |         |                |              |
    +--------+--------+--------+--------------+
                      |
                      v
       Yandex LLM Generation (temperature=0.0)
                      |
                      v
          Structured Clinical Response
\end{verbatim}

\subsection{Key Configuration
Parameters}\label{key-configuration-parameters}

\begin{longtable}[]{@{}
  >{\raggedright\arraybackslash}p{(\linewidth - 4\tabcolsep) * \real{0.3333}}
  >{\raggedright\arraybackslash}p{(\linewidth - 4\tabcolsep) * \real{0.3333}}
  >{\raggedright\arraybackslash}p{(\linewidth - 4\tabcolsep) * \real{0.3333}}@{}}
\toprule\noalign{}
\begin{minipage}[b]{\linewidth}\raggedright
Parameter
\end{minipage} & \begin{minipage}[b]{\linewidth}\raggedright
Value
\end{minipage} & \begin{minipage}[b]{\linewidth}\raggedright
Source
\end{minipage} \\
\midrule\noalign{}
\endhead
\bottomrule\noalign{}
\endlastfoot
Embedding model (documents) & \texttt{text-search-doc/latest} & Yandex
Foundation Models \\
Embedding model (queries) & \texttt{text-search-query/latest} & Yandex
Foundation Models \\
Embedding dimensionality & 256 & Yandex API \\
Retrieval top-K & 5 & K × L2 grid sweep \\
Max L2 distance threshold & 1.2 & K × L2 grid sweep \\
Chunk size & 400 words & Chunk-size grid sweep \\
Chunk overlap & 30 words & Fixed \\
Generation model & YandexGPT/latest & Yandex Foundation Models \\
Generation temperature & 0.0 & Fixed (deterministic) \\
Routing model & YandexGPT/latest & Yandex Foundation Models \\
\end{longtable}

The current registry includes four specialty agents: cardiology,
endocrinology, gastroenterology, and infectious diseases (the
``infectionist'' agent is the internal designation for the
infectious-disease specialty). Each agent owns its own FAISS dense and
BM25 sparse index; adding a further specialty is a single registry entry
plus a corpus and FAISS+BM25 index build, and every per-domain loop in
the evaluation pipeline iterates over the registered specialties rather
than hardcoded 2-specialty tuples.

\subsection{Design Decisions}\label{design-decisions}

\begin{itemize}
\tightlist
\item
  \textbf{Keyword stripping:} Raw chunk files contain a
  \texttt{KEYWORDS:} header line (produced by TF-IDF extraction). This
  line is stripped from \texttt{page\_content} before embedding to
  prevent semantic distortion. Keywords are preserved in document
  \texttt{metadata} for potential future hybrid search.
\item
  \textbf{Asymmetric embeddings:} Yandex provides separate document and
  query embedding models, optimized for asymmetric retrieval (short
  query vs.~long passage).
\item
  \textbf{Strict generation prompting:} Agent system prompts explicitly
  prohibit introducing medical facts not present in the retrieved
  context. An ``Insufficient evidence'' fallback instruction is placed
  at the end of the prompt --- a heuristic intended to exploit the LLM's
  tendency to weight later prompt content more heavily.
\end{itemize}

\subsection{3.1 Retrieval Confidence
Interpretation}\label{retrieval-confidence-interpretation}

FAISS returns an L2 distance metric for each retrieved chunk (lower is
better, max threshold = 1.2). To make this interpretable for end users,
the L2 distance is converted into a percentage confidence score using
the formula \texttt{sim\ =\ max(0,\ 1\ -\ L2\ /\ MAX\_L2\_DISTANCE)}.
This score is now displayed alongside the retrieved evidence in the
generated output.

Across the full 200-case golden dataset (50 cases per specialty × K=5 =
\textbf{1,000 chunks}, all within \texttt{MAX\_L2\_DISTANCE\ =\ 1.2}),
the confidence scores demonstrate a consistent operational range:

\begin{longtable}[]{@{}lrrrr@{}}
\toprule\noalign{}
Specialty & n chunks & Mean & Min & Max \\
\midrule\noalign{}
\endhead
\bottomrule\noalign{}
\endlastfoot
Cardiology & 250 & 19.3\% & 9.0\% & 30.8\% \\
Endocrinology & 250 & 22.5\% & 10.9\% & 33.6\% \\
Gastroenterology & 250 & 17.0\% & 8.9\% & 27.2\% \\
Infectiology & 250 & 12.9\% & 2.6\% & 20.8\% \\
\end{longtable}

While the absolute percentages appear mathematically low, they represent
highly relevant semantic matches within the 256-dimensional embedding
space (L2=0 is an exact string match, which never occurs for
natural-language Q\&A). The lower infectiology range is consistent with
that corpus's higher mean min-L2 (§5.5 reports μ=0.894 for infect vs
0.874 for cardio) --- i.e.~infect queries land further from their
nearest neighbours on average, not that the retrieval is worse. These
distributions provide an interpretable baseline for system monitoring
and observability.

\begin{center}\rule{0.5\linewidth}{0.5pt}\end{center}

\section{4. Knowledge Base}\label{knowledge-base}

\subsection{4.1 Source Documents}\label{source-documents}

\begin{longtable}[]{@{}
  >{\raggedright\arraybackslash}p{(\linewidth - 8\tabcolsep) * \real{0.1579}}
  >{\raggedleft\arraybackslash}p{(\linewidth - 8\tabcolsep) * \real{0.2105}}
  >{\raggedleft\arraybackslash}p{(\linewidth - 8\tabcolsep) * \real{0.2105}}
  >{\raggedleft\arraybackslash}p{(\linewidth - 8\tabcolsep) * \real{0.2105}}
  >{\raggedleft\arraybackslash}p{(\linewidth - 8\tabcolsep) * \real{0.2105}}@{}}
\toprule\noalign{}
\begin{minipage}[b]{\linewidth}\raggedright
Category
\end{minipage} & \begin{minipage}[b]{\linewidth}\raggedleft
Cardiology
\end{minipage} & \begin{minipage}[b]{\linewidth}\raggedleft
Endocrinology
\end{minipage} & \begin{minipage}[b]{\linewidth}\raggedleft
Gastroenterology
\end{minipage} & \begin{minipage}[b]{\linewidth}\raggedleft
Infectiology
\end{minipage} \\
\midrule\noalign{}
\endhead
\bottomrule\noalign{}
\endlastfoot
Articles & 7 & 67 & 232 & 25 \\
Cases & 59 & 74 & 167 & 59 \\
Guidelines & 113 & 101 & 104 & 212 \\
Handbooks & 190 & 0 & 8 & 9 \\
Textbooks & 25 & 914 & 30 & 165 \\
\textbf{Total documents} & \textbf{394} & \textbf{1,156} & \textbf{541}
& \textbf{470} \\
\end{longtable}

\textbf{Total across all four specialties: 2,561 source documents.}

\subsection{4.2 Processed Chunks (400 words, 30-word
overlap)}\label{processed-chunks-400-words-30-word-overlap}

\begin{longtable}[]{@{}
  >{\raggedright\arraybackslash}p{(\linewidth - 8\tabcolsep) * \real{0.1579}}
  >{\raggedleft\arraybackslash}p{(\linewidth - 8\tabcolsep) * \real{0.2105}}
  >{\raggedleft\arraybackslash}p{(\linewidth - 8\tabcolsep) * \real{0.2105}}
  >{\raggedleft\arraybackslash}p{(\linewidth - 8\tabcolsep) * \real{0.2105}}
  >{\raggedleft\arraybackslash}p{(\linewidth - 8\tabcolsep) * \real{0.2105}}@{}}
\toprule\noalign{}
\begin{minipage}[b]{\linewidth}\raggedright
Category
\end{minipage} & \begin{minipage}[b]{\linewidth}\raggedleft
Cardiology
\end{minipage} & \begin{minipage}[b]{\linewidth}\raggedleft
Endocrinology
\end{minipage} & \begin{minipage}[b]{\linewidth}\raggedleft
Gastroenterology
\end{minipage} & \begin{minipage}[b]{\linewidth}\raggedleft
Infectiology
\end{minipage} \\
\midrule\noalign{}
\endhead
\bottomrule\noalign{}
\endlastfoot
Articles & 10 & 2,487 & 2,817 & 303 \\
Cases & 105 & 1,046 & 1,073 & 390 \\
Guidelines & 1,478 & 6,916 & 4,391 & 3,116 \\
Handbooks & 972 & 0 & 53 & 434 \\
Textbooks & 5,165 & 27,342 & 690 & 3,469 \\
\textbf{Total chunks (post-filter)} & \textbf{7,730} & \textbf{37,791} &
\textbf{8,670} & \textbf{7,476} \\
\end{longtable}

\textbf{Total across all four FAISS indices: 61,667 chunks.} Pre-filter
raw chunk counts (9,024 gastro / 7,712 infect) are reduced by a
mean-word-length filter applied during index construction, which drops
354 gastro + 236 infect PDF-extraction-artifact chunks (concatenated
text with no inter-word spaces) before embedding; cardio and endo are
unaffected. See §7 L3 for the rationale.

A concrete failure surfaced when adding the gastroenterology and
infectious-diseases data: PDF extraction occasionally produced text
without word spacing, yielding tokens too long to embed reliably and
motivating the mean-word-length filter described in §7 L3.

\subsection{4.3 Chunk Size Optimization}\label{chunk-size-optimization}

A grid search over chunk sizes \{100, 200, 400, 500, 600\} words was
performed using the chunk-size grid sweep. Because re-chunking and
re-embedding the full 45,521-chunk corpus at each candidate size would
require \textasciitilde15 hours of Yandex Embedding API calls per size,
the search was run on a \textbf{proxy subset}: up to 10 keyword-relevant
documents per specialty (≤20 documents total, selected by overlap with
golden-dataset keywords), re-chunked in memory and embedded into
temporary FAISS indices.

On this proxy subset (n = 30 queries, \textasciitilde20 documents),
chunk size \textbf{400 words} achieved the highest Hit Rate at
\textbf{80.0\%}. The proxy Hit Rate is lower than the full-index Hit
Rate (96.7\%) because the subset contains only \textasciitilde20
documents vs.~the full corpus of \textasciitilde660 source documents ---
many golden-dataset queries match on documents outside the proxy subset.
The 80.0\% figure should be interpreted as a \textbf{relative ranking}
across chunk sizes, not an absolute performance estimate.

The winning chunk size (400 words) was then validated by building the
full production FAISS indices (7,730 + 37,791 chunks) and running the
complete evaluation suite, confirming 96.7\% Hit Rate on the full
30-query golden dataset.

\begin{longtable}[]{@{}
  >{\raggedright\arraybackslash}p{(\linewidth - 4\tabcolsep) * \real{0.3333}}
  >{\raggedright\arraybackslash}p{(\linewidth - 4\tabcolsep) * \real{0.3333}}
  >{\raggedright\arraybackslash}p{(\linewidth - 4\tabcolsep) * \real{0.3333}}@{}}
\toprule\noalign{}
\begin{minipage}[b]{\linewidth}\raggedright
Chunk Size (words)
\end{minipage} & \begin{minipage}[b]{\linewidth}\raggedright
Proxy Hit Rate
\end{minipage} & \begin{minipage}[b]{\linewidth}\raggedright
Notes
\end{minipage} \\
\midrule\noalign{}
\endhead
\bottomrule\noalign{}
\endlastfoot
100 & \textasciitilde53\% & Too short --- clinical context fragmented
across chunks \\
200 & \textasciitilde67\% & Original default; reasonable but
suboptimal \\
\textbf{400} & \textbf{80\%} & \textbf{Selected} --- best balance of
context and embedding quality \\
500 & \textasciitilde73\% & Token limit truncation begins affecting some
chunks \\
600 & \textasciitilde67\% & \textgreater2,048 tokens for many chunks;
forced truncation degrades quality \\
\end{longtable}

Note: the chunk-size grid was also run on the 30-case dev split using a
\textasciitilde20-document proxy subset (cost-saving). The selected
chunk size (400 words) was applied to the full corpus before the §5.8
held-out evaluation.

\subsection{4.4 Retrieval Hyperparameter Grid
Search}\label{retrieval-hyperparameter-grid-search}

A grid search over K ∈ \{3,5,7,10,15\} × L2 ∈
\{0.8,1.0,1.2,1.4,1.6,2.0\} was performed on the 30-case development
split (the grid was run on the cardiology + endocrinology development
split; re-running on the four-specialty split is outstanding work).
Hyperparameter selection was therefore performed on a strict subset of
the cases reported in §5; the §5.8 held-out test split (n = 140) reports
performance on cases never seen during tuning.

\begin{longtable}[]{@{}
  >{\raggedright\arraybackslash}p{(\linewidth - 12\tabcolsep) * \real{0.1429}}
  >{\raggedright\arraybackslash}p{(\linewidth - 12\tabcolsep) * \real{0.1429}}
  >{\raggedright\arraybackslash}p{(\linewidth - 12\tabcolsep) * \real{0.1429}}
  >{\raggedright\arraybackslash}p{(\linewidth - 12\tabcolsep) * \real{0.1429}}
  >{\raggedright\arraybackslash}p{(\linewidth - 12\tabcolsep) * \real{0.1429}}
  >{\raggedright\arraybackslash}p{(\linewidth - 12\tabcolsep) * \real{0.1429}}
  >{\raggedright\arraybackslash}p{(\linewidth - 12\tabcolsep) * \real{0.1429}}@{}}
\toprule\noalign{}
\begin{minipage}[b]{\linewidth}\raggedright
K
\end{minipage} & \begin{minipage}[b]{\linewidth}\raggedright
L2 ≤ 0.8
\end{minipage} & \begin{minipage}[b]{\linewidth}\raggedright
L2 ≤ 1.0
\end{minipage} & \begin{minipage}[b]{\linewidth}\raggedright
L2 ≤ 1.2
\end{minipage} & \begin{minipage}[b]{\linewidth}\raggedright
L2 ≤ 1.4
\end{minipage} & \begin{minipage}[b]{\linewidth}\raggedright
L2 ≤ 1.6
\end{minipage} & \begin{minipage}[b]{\linewidth}\raggedright
L2 ≤ 2.0
\end{minipage} \\
\midrule\noalign{}
\endhead
\bottomrule\noalign{}
\endlastfoot
3 & 0.0\% & 70.0\% & 86.7\% & 86.7\% & 86.7\% & 86.7\% \\
\textbf{5} & \textbf{0.0\%} & \textbf{73.3\%} & \textbf{96.7\%
\textless{}} & \textbf{96.7\%} & \textbf{96.7\%} & \textbf{96.7\%} \\
7 & 0.0\% & 76.7\% & 100.0\% & 100.0\% & 100.0\% & 100.0\% \\
10 & 0.0\% & 76.7\% & 100.0\% & 100.0\% & 100.0\% & 100.0\% \\
15 & 0.0\% & 76.7\% & 100.0\% & 100.0\% & 100.0\% & 100.0\% \\
\end{longtable}

\begin{figure}
\centering
\pandocbounded{\includesvg[keepaspectratio]{figures/F1_k_l2_heatmap.svg}}
\caption{Retrieval Hit Rate as a function of top-K and L2 threshold on
the development split. The chosen operating point (K=5, L2 ≤ 1.2) yields
the configured retrieval setting used in the rest of the evaluation.}
\end{figure}

\textbf{Chosen operating point: K=5, L2 ≤ 1.2 (96.7\%)}

\textbf{Justification for K=5 over K=7:} While K=7 achieves 100\% Hit
Rate, the additional 2 chunks per query increase the context window fed
to the LLM by 40\%. Wider retrieval windows dilute the gold-document
signal: a larger K brings in more peripheral chunks per query, which
increases the chance the LLM synthesises across loosely related context
and fabricates clinical details not grounded in any single retrieved
source. K=5 was selected as the operating point: near-perfect dev-set
Hit Rate (96.7\%) with a compact context window. A tighter K is expected
on theoretical grounds to reduce hallucination from peripheral chunks;
no K-vs-faithfulness ablation was run to confirm this on the test split.

Key observations from the grid: - \textbf{L2 ≤ 0.8 is too strict}: zero
hits across all K values --- no document vectors fall within this
radius. - \textbf{L2 = 1.2 is the critical threshold}: Hit Rate jumps
dramatically between L2=1.0 and L2=1.2 for all K values. -
\textbf{Beyond L2 = 1.2, performance plateaus}: no additional hits are
gained by relaxing the threshold further.

\subsection{4.5 Effect of Metadata Pollution on Embedding
Quality}\label{effect-of-metadata-pollution-on-embedding-quality}

\subsubsection{Hypothesis}\label{hypothesis}

Raw chunk files contain a \texttt{KEYWORDS:} header line (produced by
TF-IDF extraction). The hypothesis was that including these dense,
non-natural language keyword lists directly within the text chunk
distorts the semantic vector produced by the embedding model, thereby
degrading retrieval performance. The hypothesis follows from how dense
retrievers like DPR (\citeproc{ref-karpukhin2020dpr}{Karpukhin et al.,
2020}) are trained --- on natural-language passage / question pairs ---
so non-natural tokens occupy unusual regions of the embedding space and
pull the chunk's vector toward those regions.

\subsubsection{Experiment}\label{experiment}

Retrieval performance on the cardiology index was evaluated before and
after implementing a strict keyword-stripping pre-processing step
(removing the \texttt{KEYWORDS:} line from \texttt{page\_content} before
embedding, while preserving it in document metadata).

\subsubsection{Result (original,
confounded)}\label{result-original-confounded}

Removing metadata pollution was originally reported to increase the
cardiology Hit Rate from 93.3\% to 96.7\% (an earlier comparison,
single-case improvement on \texttt{cardio\_12}). However, that
comparison switched chunk size from 200→400 words at the same time, so
the +3.4 pp cannot be attributed to keyword stripping in isolation.

\subsubsection{Result (unconfounded 2×2
ablation)}\label{result-unconfounded-22-ablation}

A 2×2 ablation on the cardiology corpus rebuilt three of the four cells
with the existing 400-word chunks reconstructed back into raw text and
re-chunked at the target size (raw cardiology documents are no longer on
disk locally, so reconstruction proceeds from the native 400-word
chunks; documented methodological caveat). Evaluated on the 30-case
development split's cardiology slice (15 cases; 14 contributing to
Recall@5). Cells:

\begin{longtable}[]{@{}
  >{\raggedright\arraybackslash}p{(\linewidth - 12\tabcolsep) * \real{0.1429}}
  >{\raggedright\arraybackslash}p{(\linewidth - 12\tabcolsep) * \real{0.1429}}
  >{\raggedright\arraybackslash}p{(\linewidth - 12\tabcolsep) * \real{0.1429}}
  >{\raggedright\arraybackslash}p{(\linewidth - 12\tabcolsep) * \real{0.1429}}
  >{\raggedright\arraybackslash}p{(\linewidth - 12\tabcolsep) * \real{0.1429}}
  >{\raggedright\arraybackslash}p{(\linewidth - 12\tabcolsep) * \real{0.1429}}
  >{\raggedright\arraybackslash}p{(\linewidth - 12\tabcolsep) * \real{0.1429}}@{}}
\toprule\noalign{}
\begin{minipage}[b]{\linewidth}\raggedright
Cell
\end{minipage} & \begin{minipage}[b]{\linewidth}\raggedright
Chunk size
\end{minipage} & \begin{minipage}[b]{\linewidth}\raggedright
Keywords
\end{minipage} & \begin{minipage}[b]{\linewidth}\raggedright
KeywordHitRate
\end{minipage} & \begin{minipage}[b]{\linewidth}\raggedright
Recall@5
\end{minipage} & \begin{minipage}[b]{\linewidth}\raggedright
MRR@5
\end{minipage} & \begin{minipage}[b]{\linewidth}\raggedright
n
\end{minipage} \\
\midrule\noalign{}
\endhead
\bottomrule\noalign{}
\endlastfoot
A (historical) & 200 & keep & 93.3\% on cardio\_1..30 (different case
set + older code path; not directly comparable) & --- & --- & --- \\
\textbf{B} & 200 & strip & \textbf{100.0\%} (15/15) & \textbf{59.5\%} &
0.893 & 14 \\
\textbf{C} & 400 & keep & \textbf{93.3\%} (14/15) & \textbf{69.0\%} &
0.881 & 14 \\
\textbf{D} (current production) & 400 & strip & \textbf{93.3\%} (14/15)
& \textbf{69.0\%} & 0.875 & 14 \\
\end{longtable}

\textbf{Decomposition (on the comparable B/C/D cells):}

\begin{itemize}
\tightlist
\item
  \textbf{Main effect of keyword stripping (at chunk\_size = 400, i.e.~D
  − C):} 0.0 pp on KeywordHitRate, 0.0 pp on Recall@5. The two cells are
  identical on every grouped metric. Keyword stripping on its own
  contributes \textbf{nothing measurable} at 400-word chunk size on the
  dev cardiology slice.
\item
  \textbf{Main effect of chunk size (at strip = True, i.e.~D − B):} −6.7
  pp on KeywordHitRate (200 wins: 100.0\% vs 93.3\%), +9.5 pp on
  Recall@5 (400 wins: 69.0\% vs 59.5\%). The two metrics point opposite
  directions because 200-word chunks fragment each document into
  \textasciitilde2× more pieces --- more chunks means more chances for
  any expected keyword to appear in the top-5 (boosting KeywordHitRate),
  but top-5 then covers fewer unique source documents (depressing the
  doc-level Recall@5).
\item
  \textbf{Interaction (D − C) − (B − A):} with A unmeasured on the
  current dev split (raw 200-word chunks not on disk, A's historical
  93.3\% was on the old \texttt{cardio\_1..30} 30-case set), the
  interaction term cannot be computed strictly. Imputing A ≈ B = 100\%
  (since the strip effect at 400 is 0, the strip effect at 200 is
  expected to be near 0 too): interaction ≈ (93.3 − 93.3) − (100 − 100)
  = \textbf{0.0 pp}.
\end{itemize}

\textbf{Corrected claim.} The +3.4 pp originally attributed to keyword
stripping in the prior confounded comparison (``Hit Rate improved from
93.3\% to 96.7\%'') is \textbf{0.0 pp from keyword stripping}, plus a
sample-size-dependent chunk-size effect that flips sign depending on
which metric is read. The original narrative confused a one-case
(cardio\_12) improvement, which on n = 30 was +3.3 pp, with a real
effect of stripping --- but on the current dev split with chunk size
held constant the strip toggle produces a 0-case difference. With Wilson
95\% CIs on the dev split's small n (15 cardio cases), a 1-case swing is
±7 pp noise; the historical +3.3 pp is well inside that noise band. The
chunk-size effect is also small in absolute terms (≤ 1 case on Recall@5
differences) and metric-dependent --- neither chunk size nor keyword
stripping moves cardiology retrieval quality by a margin that survives
the noise band on this dev split. What does matter, on a much larger
scale, is choice of retriever (dense vs sparse vs hybrid) --- see §5.3.2
for the BM25 ablation, which shows FAISS beating BM25 by 22.2 pp on
Recall@5, much larger than any chunk-size / strip effect documented
here. (This ablation was run on the cardiology corpus only;
generalisation to the other three specialties is unverified.)

\begin{center}\rule{0.5\linewidth}{0.5pt}\end{center}

\section{5. Evaluation}\label{evaluation}

All evaluations use a \textbf{golden dataset} of 200 clinical cases
across three difficulty tiers (Core, Peripheral, Out-of-Scope) and four
specialties (cardiologist, endocrinologist, gastroenterologist,
infectionist --- 50 cases per specialty). The dev/test split is
\texttt{case-number\ ≤\ 15} → dev (60 cases) / rest → test (140 cases).
This 4-specialty evaluation widens every number in §5 from the prior
2-specialist scope; the cardiology and endocrinology test cases are
byte-identical to the prior 2-specialist baseline, so any cardio / endo
cell that retains its prior value is a verified regression check (per
§5.10 below).

\subsection{5.1 Routing Architecture: LLM vs.~Keyword vs.~TF-IDF
Baselines}\label{routing-architecture-llm-vs.-keyword-vs.-tf-idf-baselines}

Two non-LLM baselines are reported alongside the LLM router on the
\textbf{held-out test split (n = 140)}: a 4-specialty keyword dictionary
(the keyword-route baseline implementation extended to four specialties
--- adds GASTRO\_KEYWORDS + INFECT\_KEYWORDS sets with a ``highest hit
count wins, registry-order tie-break'' rule) and a TF-IDF (1--2 grams) +
LogisticRegression model trained on the development split.

\begin{longtable}[]{@{}
  >{\raggedright\arraybackslash}p{(\linewidth - 10\tabcolsep) * \real{0.1667}}
  >{\raggedright\arraybackslash}p{(\linewidth - 10\tabcolsep) * \real{0.1667}}
  >{\raggedright\arraybackslash}p{(\linewidth - 10\tabcolsep) * \real{0.1667}}
  >{\raggedright\arraybackslash}p{(\linewidth - 10\tabcolsep) * \real{0.1667}}
  >{\raggedright\arraybackslash}p{(\linewidth - 10\tabcolsep) * \real{0.1667}}
  >{\raggedright\arraybackslash}p{(\linewidth - 10\tabcolsep) * \real{0.1667}}@{}}
\toprule\noalign{}
\begin{minipage}[b]{\linewidth}\raggedright
Method
\end{minipage} & \begin{minipage}[b]{\linewidth}\raggedright
Cardiology
\end{minipage} & \begin{minipage}[b]{\linewidth}\raggedright
Endocrinology
\end{minipage} & \begin{minipage}[b]{\linewidth}\raggedright
Gastroenterology
\end{minipage} & \begin{minipage}[b]{\linewidth}\raggedright
Infectiology
\end{minipage} & \begin{minipage}[b]{\linewidth}\raggedright
Overall
\end{minipage} \\
\midrule\noalign{}
\endhead
\bottomrule\noalign{}
\endlastfoot
Keyword Baseline & 94.3\% (33/35) {[}81.4\%--98.4\%{]} & 97.1\% (34/35)
{[}85.5\%--99.5\%{]} & 80.0\% (28/35) {[}64.1\%--90.0\%{]} & 45.7\%
(16/35) {[}30.5\%--61.8\%{]} & 79.3\% (111/140) {[}71.8\%--85.2\%{]} \\
TF-IDF Baseline (4-spec dev-trained) & 54.3\% (19/35)
{[}38.2\%--69.5\%{]} & 71.4\% (25/35) {[}54.9\%--83.7\%{]} & 74.3\%
(26/35) {[}57.9\%--85.8\%{]} & 80.0\% (28/35) {[}64.1\%--90.0\%{]} &
70.0\% (98/140) {[}62.0\%--77.0\%{]} \\
LLM Router & 100.0\% (35/35) {[}90.1\%--100.0\%{]} & 97.1\% (34/35)
{[}85.5\%--99.5\%{]} & 88.6\% (31/35) {[}74.0\%--95.5\%{]} & 97.1\%
(34/35) {[}85.5\%--99.5\%{]} & \textbf{95.7\% (134/140)
{[}91.0\%--98.0\%{]}} \\
\end{longtable}

\emph{(Wilson 95\% CIs via \texttt{statsmodels}. Test split n = 140 = 35
cases per specialty (cases 16--50). The TF-IDF model was retrained on
the 60-case 4-specialty dev split on 2026-05-22 with
\texttt{random\_state=42}; the model now carries all four
\texttt{classes\_}. Pre-2026-05-22 it was 2-class (cardio + endo only)
and reported 0.0\% on gastro/infect --- see §7 L9 for the historical
context.)}

\textbf{Interpretation.} The LLM Router's overall 95.7\% (134/140,
Wilson lower 91.0\%) is \textbf{16.4 pp above the 4-class keyword
baseline (79.3\%)} and \textbf{25.7 pp above the dev-trained TF-IDF
baseline (70.0\%)} --- the post-retrain TF-IDF gap is meaningfully
tighter than the pre-retrain 56.4 pp.~Cardiology routing is 35/35 =
100\% --- identical to the prior 2-specialty baseline, a verified
regression check. The 6 LLM misses are all defensible cross-specialty
ambiguities: \texttt{endo\_48} (SGLT2-inhibitor complication → UTI/DKA)
routed to infectionist; \texttt{gastro\_17} (NAFLD in metabolic
syndrome) and \texttt{gastro\_22} (Wilson's disease with
tremor/dysarthria) routed to endocrinologist; \texttt{gastro\_37}
(post-buffet diarrhoea+vomiting+fever) routed to infectionist;
\texttt{gastro\_39} (haemochromatosis presenting with diabetes +
arthralgia) routed to endocrinologist; \texttt{infect\_17} (chronic HCV)
routed to gastroenterologist. Each of these is a clinically reasonable
alternate routing where the case spans two specialties; they are not
bugs.

The keyword baseline's catastrophic 45.7\% on the infectiology column is
the consequence of a historical endocrinology fallback rule --- when no
keyword set matches, the rule returns \texttt{endocrinologist} --- and
many T1 infect cases (CMV reactivation post-transplant, febrile
neutropenia, catheter-associated UTI, etc.) use surface vocabulary that
fails the infect-specific keyword set but matches an endocrine fallback
class. The retrained TF-IDF reverses this pattern: it is
\emph{strongest} on infectiology (80.0\%) and gastroenterology (74.3\%)
--- where the new corpora have distinctive vocabulary --- and
\emph{weakest} on cardiology (54.3\%), where T2/T3 cases share more
surface tokens with endocrinology (\texttt{hypertension},
\texttt{diabetes}, \texttt{metabolic}). \textbf{The case for the LLM is
now made on the clear-domain test set itself}, not just on ambiguous
cases as in the prior 2-specialty configuration: at 4-specialty scope
the LLM clears the \emph{stronger} of the two non-LLM baselines (keyword
at 79.3\%) by 16.4 pp, well outside any CI overlap.

\subsection{5.2 Qualitative Behaviour on Cross-Domain Ambiguous
Cases}\label{qualitative-behaviour-on-cross-domain-ambiguous-cases}

To probe the router's behaviour on clinically ambiguous queries, a
dedicated test set of \textbf{14 cross-domain cases} was constructed: 8
cardio×endo cases from the original ambiguous-case set
(\texttt{ambig\_1..8}) plus 6 cases spanning the new specialty pairs
added with the four-specialty expansion (\texttt{ambig\_9..14}). Each
case intentionally spans two specialties --- no single routing decision
is strictly ``correct''; the LLM router's observed routing decision is
compared against the keyword and 4-class TF-IDF baseline predictions
(after the 2026-05-22 retrain --- see §5.1, §7 L9 --- both baselines can
now return any of the four specialties).

\textbf{LLM stays inside \texttt{valid\_domains} on all 14 ambiguous
cases}, while the keyword baseline misroutes \texttt{ambig\_2..ambig\_8}
to cardiologist because cardiology keywords dominate by surface count.
Appendix A.1 lists the full 14-case breakdown with each router's
prediction. The LLM routes \texttt{ambig\_8} (carcinoid heart disease)
to \textbf{gastroenterologist}, which has clinical merit --- carcinoid
tumours are gastrointestinal in origin (mid-gut neuroendocrine) and the
syndrome is treated by removing the GI primary, not by valve surgery
alone. \texttt{ambig\_8} was originally annotated before the
four-specialty expansion with
\texttt{{[}cardiologist,\ endocrinologist{]}}; widened on 2026-05-22 to
include \texttt{gastroenterologist}, so the LLM's choice now counts as
inside \texttt{valid\_domains} (the same 2026-05-22 pass unified the
schema on \texttt{valid\_domains} and retired the legacy
\texttt{domains} field, now resolved).

Across the 14 ambiguous cases the \textbf{LLM produces a 5+3 cardio/endo
split on \texttt{ambig\_1..7}} (\texttt{ambig\_8} lands in gastro after
the four-specialty expansion as discussed above) and routes the 6 new
cases with clinical intuition: \texttt{ambig\_9} to gastro (the
peptic-ulcer side has actionable management), \texttt{ambig\_10} to endo
(the new hyperglycaemia is what's outside the trauma physician's normal
toolkit), \texttt{ambig\_11} to cardio (DCM on ART is a cardiologic
management problem; HIV is stable), \texttt{ambig\_12} to gastro (AIH is
the actionable diagnosis), \texttt{ambig\_13} to infect (C. difficile is
an antimicrobial-management problem), \texttt{ambig\_14} to infect (SBP
is an antimicrobial-management problem on top of a known cirrhotic). All
6 land within their stated \texttt{valid\_domains}. The keyword baseline
now produces a more varied distribution because the 4-specialty
extension lets it pick gastro or infect when one of those keyword sets
fires more often (\texttt{ambig\_9} → gastro, \texttt{ambig\_10} → endo,
\texttt{ambig\_11} → infect, \texttt{ambig\_12} → endo,
\texttt{ambig\_13} and \texttt{ambig\_14} → gastro), but it still
misroutes \texttt{ambig\_2..ambig\_8} to cardiologist because cardiology
keywords dominate by surface count.

\subsubsection{5.2.1 Adversarial Routing}\label{adversarial-routing}

The §5.1 / §5.2 numbers above are computed on the 140-case held-out test
split, whose queries are written in clean English with standard medical
vocabulary. To stress-test the router's robustness on inputs that
violate those assumptions, a dedicated 64-case adversarial test set
(tier=4 / \texttt{tier\_label="adversarial"}) was constructed across
four categories --- each with 16 cases (8 cardio×endo from the original
adversarial set + 8 gastro/infect added with the four-specialty
expansion). Per-category accuracy with Wilson 95\% CIs from the LLM
router:

\begin{longtable}[]{@{}llll@{}}
\toprule\noalign{}
Category & n & Correct & Accuracy {[}Wilson 95\% CI{]} \\
\midrule\noalign{}
\endhead
\bottomrule\noalign{}
\endlastfoot
\texttt{misspelled} & 16 & 16 & 100.0\% {[}80.6\%--100.0\%{]} \\
\texttt{non\_english} & 16 & 16 & 100.0\% {[}80.6\%--100.0\%{]} \\
\texttt{dominant\_pathology\_mismatch} & 16 & 16 & 100.0\%
{[}80.6\%--100.0\%{]} \\
\texttt{symptom\_only\_ambiguous} & 16 & 15 & 93.8\%
{[}71.7\%--98.9\%{]} \\
\textbf{Overall} & \textbf{64} & \textbf{63} & \textbf{98.4\%
{[}91.7\%--99.7\%{]}} \\
\end{longtable}

\emph{(\texttt{symptom\_only\_ambiguous} cases carry a
\texttt{valid\_domains} field listing the two specialties either of
which is counted correct because the query names only symptoms. The
other three categories use a single \texttt{expected\_specialist} set by
the test author. The doubling from 32 → 64 cases tightens every Wilson
lower bound by \textasciitilde13 pp --- from 67.6\% to 80.6\% per
category --- while the overall lower bound tightens from 89.3\% to
91.7\%.)}

The adversarial headline of 98.4\% (63/64, Wilson lower 91.7\%) is only
2.7 pp below the clear-domain test headline of 95.7\% (134/140) ---
i.e.~the router is consistent across written-clean-English and
stress-test inputs alike. Three notable observations:

\begin{itemize}
\tightlist
\item
  The \texttt{misspelled} and \texttt{non\_english} categories remain at
  16/16 = 100.0\% after doubling --- \texttt{misspelled} queries with
  typos like \texttt{tubercolosis}, \texttt{feber}, \texttt{lypase},
  \texttt{endscopy} and Cyrillic/Spanish/French queries with the
  gastro/infect surface vocabulary all route correctly. YandexGPT's
  Russian-native training is the obvious driver for the Russian queries;
  the French/Spanish performance is consistent with the model's
  behaviour on the original 2-spec adversarial set.
\item
  The single \texttt{symptom\_only\_ambiguous} miss is
  \texttt{adv\_amb\_6} (from the original adversarial cohort: 70-y/o
  with weight loss, reduced appetite, and generalised weakness with no
  focal symptoms). The LLM routed to \texttt{gastroenterologist},
  outside the case's
  \texttt{valid\_domains\ =\ {[}cardiologist,\ endocrinologist{]}}. The
  case's \texttt{valid\_domains} was authored before gastroenterology
  was registered; routing to gastro is a clinically reasonable response
  to the geriatric weight-loss presentation (occult GI malignancy
  work-up) --- same provenance as \texttt{ambig\_8} (carcinoid →
  gastro). Widening \texttt{adv\_amb\_6.valid\_domains} to include
  \texttt{gastroenterologist} is the dataset-curation fix. The dataset
  has been updated, so the next routing eval pass will count this case
  as a pass (lifting \texttt{symptom\_only\_ambiguous} to 16/16 =
  100.0\% and the adversarial overall to 64/64 = 100.0\%); the table
  above still reflects the historical 63/64 eval run.
\item
  The \texttt{dominant\_pathology\_mismatch} category was the hardest by
  design (surface vocabulary deliberately points to the opposite
  specialty from the actionable diagnosis) and stays at 16/16 = 100.0\%.
  The 8 new cases here include ``Crohn's on infliximab develops cavitary
  TB'' (→ infectionist, not gastroenterologist), ``post-CAP develops
  heart failure'' (→ cardiologist, not infectionist), and ``long-term
  PPI → hypocalcaemic tetany'' (→ endocrinologist, not
  gastroenterologist) --- all routed correctly despite the misleading
  surface vocabulary.
\end{itemize}

\subsection{5.3 Retrieval Hit Rate}\label{retrieval-hit-rate}

Each query is sent to the correct specialist agent (bypassing the
router). The agent retrieves K=5 chunks with L2 ≤ 1.2. \textbf{Recall@5}
is the primary grounded metric: every Tier 1/2 case carries a
\texttt{gold\_sources} annotation listing 1--3 source documents that
contain the correct answer (auto-annotated by selecting up to three
documents per case from the top-20 retrieval window with ≥1
expected-keyword hit; see the denominator note below and §7 L6 for the
circularity caveat); Recall@5 is the fraction of those gold documents
that appear among the K=5 retrieved chunks (pooled across cases ---
every gold-doc slot is one Bernoulli trial). \textbf{MRR@5} is the
reciprocal rank of the first retrieved gold document, averaged across
annotated cases. \textbf{Refusal Rate (T3)} is the fraction of Tier 3
cases where the agent's L2 distance gate fires (numeric refusal gate;
§5.5). \textbf{KeywordHitRate} is the original keyword-co-occurrence
metric kept here as a loose secondary signal for comparison.

\begin{longtable}[]{@{}
  >{\raggedright\arraybackslash}p{(\linewidth - 8\tabcolsep) * \real{0.2000}}
  >{\raggedright\arraybackslash}p{(\linewidth - 8\tabcolsep) * \real{0.2000}}
  >{\raggedright\arraybackslash}p{(\linewidth - 8\tabcolsep) * \real{0.2000}}
  >{\raggedright\arraybackslash}p{(\linewidth - 8\tabcolsep) * \real{0.2000}}
  >{\raggedright\arraybackslash}p{(\linewidth - 8\tabcolsep) * \real{0.2000}}@{}}
\toprule\noalign{}
\begin{minipage}[b]{\linewidth}\raggedright
Domain
\end{minipage} & \begin{minipage}[b]{\linewidth}\raggedright
Recall@5
\end{minipage} & \begin{minipage}[b]{\linewidth}\raggedright
MRR@5 {[}95\% CI{]}
\end{minipage} & \begin{minipage}[b]{\linewidth}\raggedright
KeywordHitRate (legacy)
\end{minipage} & \begin{minipage}[b]{\linewidth}\raggedright
Refusal Rate (T3, numeric gate)
\end{minipage} \\
\midrule\noalign{}
\endhead
\bottomrule\noalign{}
\endlastfoot
Cardiology & 56.6\% (43/76) {[}45.4\%--67.1\%{]} & 0.651
{[}0.490--0.812{]} & 89.3\% (25/28) {[}72.8\%--96.3\%{]} & 0.0\% (0/8)
{[}0.0\%--32.4\%{]} \\
Endocrinology & 55.8\% (43/77) {[}44.9\%--66.5\%{]} & 0.717
{[}0.572--0.853{]} & 92.9\% (26/28) {[}77.4\%--98.0\%{]} & 0.0\% (0/7)
{[}0.0\%--35.4\%{]} \\
Gastroenterology & 56.3\% (40/71) {[}44.8\%--67.3\%{]} & 0.731
{[}0.590--0.865{]} & 78.6\% (22/28) {[}60.5\%--89.8\%{]} & 42.9\% (3/7)
{[}15.8\%--75.0\%{]} \\
Infectiology & 56.3\% (36/64) {[}44.1\%--67.8\%{]} & 0.692
{[}0.529--0.841{]} & 78.6\% (22/28) {[}60.5\%--89.8\%{]} & 42.9\% (3/7)
{[}15.8\%--75.0\%{]} \\
\textbf{Overall} & \textbf{56.2\% (162/288) {[}50.5\%--61.9\%{]}} &
\textbf{0.698 {[}0.624--0.770{]}} & \textbf{84.6\% (95/112)
{[}76.8\%--90.2\%{]}} & \textbf{20.7\% (6/29) {[}9.8\%--38.4\%{]}} \\
\end{longtable}

\emph{(Recall@5 Wilson 95\% CIs are on the pooled gold-doc Bernoulli
(each Tier 1/2 case contributes 1--3 gold-doc trials; totals: cardio 76,
endo 77, gastro 71, infect 64 --- 288 trials across 105 annotated
test-set cases). MRR@5 95\% CIs are percentile-method bootstrap
intervals over the per-case reciprocal-rank vector (B=10000 resamples,
RNG seed=12345). The Refusal Rate column uses the numeric gate
\texttt{L2\_REJECT\_MIN\ =\ 1.020} (see §5.5 below) rather than the
legacy zero-chunk metric.)}

105 of the 111 Tier 1/2 test cases were auto-annotated by the
gold-source annotator. The remaining 6 are corpus-gap cases listed in
§5.3.1. (After the 2026-05-22 keyword-set fix for \texttt{gastro\_37}
and \texttt{infect\_14} --- see §5.3.1 retirement note --- both have
gold; \texttt{infect\_14} is in the dev split so does not affect the
test-set count, while \texttt{gastro\_37} lifts the test-set annotated
count from 104 to 105.) The 29 Tier 3 cases have
\texttt{gold\_sources:\ {[}{]}} by design and contribute only to the
Refusal Rate column. See §7 L6 for the auto-annotator circularity
caveat.

Recall@5 (56.2\%) falls far below KeywordHitRate (84.6\%) because the
two metrics measure different things. KeywordHitRate counts a case as a
hit if any of the 5 retrieved chunks contains any expected keyword
anywhere --- including adjacent, off-topic content that happens to share
a common word. Recall@5 is far stricter: it requires the \emph{specific
documents} containing the answer (annotated via top-20 keyword coverage,
then capped at 3) to land in the \emph{top-5} retrieval window. The
28-point gap is the part of the corpus that ranks 6--20 in retrieval
order --- relevant, but not surfaced at K=5.

The re-tune for the 4-specialty corpus (§5.5) lowered the per-specialty
gate strength on Tier 3 cases: the gate now fires on \textbf{20.7\%
(6/29) of T3 cases} across all four specialties --- way below the prior
target of ≥80\%. The split is uneven: cardio and endo T3 cases never
trigger the gate (0/8 and 0/7) because their min-L2 sits below the new
threshold of 1.020; gastro and infect T3 cases trigger 3/7 each (43\%),
because their corpus-mean L2 (gastro μ=0.880, infect μ=0.894) is closer
to the threshold. The trade-off is documented in §5.5; the
alternate-strategy follow-up is in §7 L7.

\subsubsection{5.3.1 Tier 2 Corpus Coverage
Audit}\label{tier-2-corpus-coverage-audit}

By design, Tier 2 (peripheral) queries stress-test the boundaries of the
knowledge base. After the 2026-05-22 retirement pass (see note below),
the audit surfaces \textbf{six} cases with empty \texttt{gold\_sources}
--- i.e.~the gold-source annotator could not find any document in the
top-20 retrieval window with ≥1 expected-keyword match. One is a
pre-existing cardio case; one is a pre-existing endo case; four are new
gastro / infect cases added with the four-specialty expansion. For all
six, retrieval surfaces \emph{adjacent} content (chunks within
\texttt{MAX\_L2\_DISTANCE\ =\ 1.2}) but the corpus does not contain a
document that matches the specific clinical-scenario keywords.

\begin{quote}
\textbf{Retirement note (2026-05-22).} Four cases that previously
appeared in this audit have been removed: - \textbf{\texttt{gastro\_37}
(acute viral gastroenteritis)} and \textbf{\texttt{infect\_14} (C.
difficile colitis)} were keyword-set mismatches, not corpus gaps.
Replacing their \texttt{expected\_keywords} with corpus-aligned
vocabulary
(\texttt{gastroenteritis}/\texttt{viral}/\texttt{dehydration}/\texttt{rotavirus}/\texttt{supportive}
for gastro\_37;
\texttt{Clostridium\ difficile}/\texttt{C.\ difficile}/\texttt{CDI}/\texttt{vancomycin}/\texttt{metronidazole}
for infect\_14) let the auto-annotator pick 3 gold docs each. Neither
lands its gold in the top-5 retrieval window (gastro T1 Recall@5 dropped
4.8 pp to 52.8\%; infect\_14 is in the dev split, so no test-split
shift). - \textbf{\texttt{cardio\_23} (pericardial effusion with
tamponade risk)} and \textbf{\texttt{cardio\_25} (Dressler syndrome /
post-pericardiotomy)} have non-empty \texttt{gold\_sources} from a prior
annotation pass --- they are not currently empty cases, so they are no
longer listed here. The underlying procedural-cardiology coverage gap
they originally illustrated persists; it is now captured solely under
the \texttt{cardio\_35} bullet and §7 L2.
\end{quote}

\textbf{Cardiology (1 case, original baseline --- preserved):} -
\textbf{\texttt{cardio\_35} (STEMI complicated by complete heart
block):} retrieval returns standard revascularisation chunks; corpus has
no electrophysiology chunks on \texttt{temporary\ pacing} /
\texttt{pacemaker} for acute blocks. Fix: electrophysiology guidelines.

\textbf{Endocrinology (1 case, original baseline --- preserved):} -
\textbf{\texttt{endo\_46} (Hypoglycaemia unawareness in long-standing
T1DM):} retrieval returns glycaemic-control chunks; corpus has no
\texttt{hypoglycaemia\ unawareness} / \texttt{closed-loop} /
\texttt{rtCGM} chunks for the specific impaired-awareness presentation.

\textbf{Gastroenterology (2 new cases, four-specialty expansion):} -
\textbf{\texttt{gastro\_39} (Hereditary haemochromatosis):} only 12
chunks match \texttt{haemochromatosis\textbar{}hemochromatosis} in the
gastro corpus; the specific keyword set (\texttt{HFE}, \texttt{C282Y},
\texttt{transferrin\ saturation}, \texttt{phlebotomy}) only partially
appears. Fix: add metabolic-liver-disease references with explicit
gene/genetics coverage. - \textbf{\texttt{gastro\_44} (Zollinger-Ellison
syndrome):} only 7 chunks match in the gastro corpus. Fix: add
neuroendocrine-tumour references with explicit gastrinoma coverage.

\textbf{Infectiology (2 new cases, four-specialty expansion):} -
\textbf{\texttt{infect\_21} (HSV encephalitis):} 139 chunks match
generic encephalitis terms but the HSV+aciclovir+temporal-lobe+CSF-PCR
vocabulary doesn't co-occur. Fix: add neurology / infectious-disease
overlap content on viral encephalitis management. -
\textbf{\texttt{infect\_39} (Prosthetic joint infection):} only 6 chunks
match \texttt{PJI}; top hit retrieves at L2 = 1.107, near the
\texttt{MAX\_L2\_DISTANCE\ =\ 1.2} ceiling. Genuine corpus coverage gap.
Fix: add orthopaedic-infectious-disease references.

Adding the underlying source material flagged above is the concrete next
ingestion step.

\subsubsection{5.3.2 Retriever Comparison: FAISS Dense vs BM25 Sparse vs
Random Baseline vs
Oracle}\label{retriever-comparison-faiss-dense-vs-bm25-sparse-vs-random-baseline-vs-oracle}

To answer the missing sparse-vs-dense baseline question flagged in §2.2,
a BM25 index (\texttt{rank-bm25==0.2.2}, lowercase / alphanumeric /
≥2-char tokens) was built over the \textbf{same chunks already indexed
by FAISS} (loaded directly from the FAISS docstore so the corpora are
identical). The table below is on the held-out 140-case test split,
pooled gold-doc Bernoulli with Wilson 95\% CI on Recall@5; MRR@5 is the
mean reciprocal rank of the first retrieved gold document.

\textbf{Dense FAISS retrieval outperforms BM25 by 22.2 percentage points
on Recall@5 overall (56.2\% vs 34.0\%) across all four specialties.}
Tested with McNemar's exact test on the 288 gold-source paired binary
outcomes, the 22.2 pp Recall@5 advantage of FAISS over BM25 is
statistically significant (discordant pairs b = 95 / c = 31, p
\textless{} 10⁻⁸; equivalently McNemar χ²(1) = 31.5 with continuity
correction, p ≈ 2.0 × 10⁻⁸). The gap is widest on the original
2-specialty conditions --- Endocrinology T1 reaches 45.4 pp (60.6\% vs
15.2\%) and Cardiology T1 reaches 33.4 pp (59.0\% vs 25.6\%) --- and
narrows on the new specialties' core conditions: Gastro T1 is 11.1 pp
(52.8\% vs 41.7\%) and Infect T1 is 18.8 pp (50.0\% vs 31.2\%). BM25
narrows the gap further on Tier 2 (cardio T2: 10.9 pp, gastro T2: 22.9
pp, infect T2: 18.7 pp) where exact-entity-name retrieval helps
(\texttt{dressler}, \texttt{colchicine}, \texttt{pericardiocentesis} in
cardio T2; \texttt{Wilson\ disease}, \texttt{Zollinger-Ellison} in
gastro T2; \texttt{Lemierre}, \texttt{mucormycosis},
\texttt{Strongyloides} in infect T2/T3). \textbf{BM25 never overtakes
FAISS on any tier of any specialty.} Random retrieval stays at
0.0--13.5\% --- a sanity floor. The four-specialty mean (56.2\%) is
within 0 pp of the prior 2-specialty mean (56.2\%, n = 153) --- i.e.~the
four-specialty extension preserves overall retrieval headline within
noise; per-cell verification in §5.10 confirms cardio/endo cells are
byte-identical.

\begin{figure}
\centering
\pandocbounded{\includesvg[keepaspectratio]{figures/F3_retriever_bar.svg}}
\caption{Recall@5 of FAISS dense retrieval vs BM25 sparse, random
baseline, and oracle ceiling, stratified by Tier (T1 / T2) and
specialty. FAISS leads BM25 by 22.2 pp overall, matching the value
reported in the §5.3.2 body prose.}
\end{figure}

Appendix A.2 contains the full per-Tier per-specialty breakdown of
Recall@5 across the four retrievers, with FAISS and BM25 MRR@5 95\%
confidence intervals.

\subsection{5.4 Faithfulness (Generation
Quality)}\label{faithfulness-generation-quality}

The full RAG pipeline (retrieval → LLM generation) is evaluated by three
independent LLM-as-a-judge models, each given the identical strict
faithfulness prompt. The methodology follows RAGAS-style automated RAG
evaluation (\citeproc{ref-es2023ragas}{Es et al., 2023}) extended with
the multi-judge minimum-rate protocol motivated by Zheng et al.'s
characterisation of same-family judge self-preference bias
(\citeproc{ref-zheng2023mtbench}{Zheng et al., 2023}). The primary judge
is YandexGPT (same family as the generation model); the secondary judge
is YandexGPT-Lite (distinct family, smaller); the tertiary judge
(configured 2026-05-23) is OpenAI's open-weights GPT-OSS-120B served via
OpenRouter, a \textbf{cross-vendor} judge that closes the residual
same-family bias of an all-Yandex panel.

\textbf{Test split (n = 140), 4-specialist evaluation, 3-judge panel:}

\begin{longtable}[]{@{}
  >{\raggedright\arraybackslash}p{(\linewidth - 12\tabcolsep) * \real{0.1429}}
  >{\raggedright\arraybackslash}p{(\linewidth - 12\tabcolsep) * \real{0.1429}}
  >{\raggedright\arraybackslash}p{(\linewidth - 12\tabcolsep) * \real{0.1429}}
  >{\raggedright\arraybackslash}p{(\linewidth - 12\tabcolsep) * \real{0.1429}}
  >{\raggedright\arraybackslash}p{(\linewidth - 12\tabcolsep) * \real{0.1429}}
  >{\raggedright\arraybackslash}p{(\linewidth - 12\tabcolsep) * \real{0.1429}}
  >{\raggedright\arraybackslash}p{(\linewidth - 12\tabcolsep) * \real{0.1429}}@{}}
\toprule\noalign{}
\begin{minipage}[b]{\linewidth}\raggedright
Judge
\end{minipage} & \begin{minipage}[b]{\linewidth}\raggedright
Provider
\end{minipage} & \begin{minipage}[b]{\linewidth}\raggedright
Model URI
\end{minipage} & \begin{minipage}[b]{\linewidth}\raggedright
Faithful
\end{minipage} & \begin{minipage}[b]{\linewidth}\raggedright
Total Judged
\end{minipage} & \begin{minipage}[b]{\linewidth}\raggedright
Rate
\end{minipage} & \begin{minipage}[b]{\linewidth}\raggedright
Wilson 95\% CI
\end{minipage} \\
\midrule\noalign{}
\endhead
\bottomrule\noalign{}
\endlastfoot
Primary & yandex & \texttt{gpt://\{folder\}/yandexgpt/latest} & 132 &
132 & 100.0\% & {[}97.2\%--100.0\%{]} \\
Secondary & yandex & \texttt{gpt://\{folder\}/yandexgpt-lite/latest} &
131 & 132 & 99.2\% & {[}95.8\%--99.9\%{]} \\
Tertiary (cross-vendor) & openrouter & \texttt{openai/gpt-oss-120b:free}
& 128 & 132 & 97.0\% & {[}92.5\%--98.8\%{]} \\
\textbf{Minimum (any-judge HALLUCINATION ⇒ HALLUCINATION)} & --- & --- &
\textbf{127} & \textbf{132} & \textbf{96.2\%} &
\textbf{{[}91.4\%--98.4\%{]}} \\
\end{longtable}

Total cases 140; 6 Tier-3 cases triggered the refusal gate (3 gastro T3
+ 3 infect T3) and are excluded from the judge denominator; one further
case (\texttt{infect\_17}) had a \texttt{None} (API-call failure)
verdict from the primary judge in the 2026-05-21 Yandex run and
additional exclusions arose from the multi-judge intersection, leaving
132 judged cases for the headline analysis (35 cardio + 35 endo + 32
gastro + 30 infect). The tertiary judge calls were made on 2026-05-23
against the cached \texttt{agent.answer(query)} re-generations of the
same 133 non-fallback cases (\texttt{MAX\_RETRIES=1}, 29.4 min
wall-clock, 0 NONE / 0 exhausted-retry failures).

Pairwise Cohen's κ on the n = 132 intersection: \textbf{κ(primary,
secondary) = 0.000 → ``degenerate (one marginal = 0; observed agreement
= 131/132)''} and \textbf{κ(primary, tertiary) = 0.000} for the same
reason --- the primary judge marks every case FAITHFUL, so its row
marginal P(HALLUCINATION) = 0 and Cohen's chance-agreement term equals
the observed-agreement term. \textbf{κ(secondary, tertiary) = −0.012} is
\emph{not} degenerate: both judges have variance, but their
HALLUCINATION verdict sets are \textbf{disjoint} (secondary flags only
\texttt{cardio\_40}; tertiary flags \texttt{cardio\_17},
\texttt{endo\_21}, \texttt{endo\_38}, \texttt{infect\_34}), so the
chance-agreement term slightly exceeds observed agreement. κ is reported
verbatim with the degeneracies flagged explicitly.

\textbf{Gwet's AC1 is the better-behaved chance-corrected statistic in
the degenerate regime.} Cohen's κ degenerates whenever one rater's
marginal class probability is 0 (i.e.~labels everything the same way),
because chance agreement under κ's assumption of independent draws from
each rater's marginal then equals the observed agreement and the
denominator vanishes. Gwet's AC1 (\citeproc{ref-gwet2008ac1}{Gwet,
2008}) replaces that with chance agreement derived from the
\textbf{empirical class prior averaged across raters} (π = total
FAITHFUL labels / 2n), which only collapses to 1.0 when both raters tie
on the same extreme --- i.e.~when there is genuinely no variance to
disagree on. AC1 is computed in the evaluation pipeline alongside κ in
the per-pair table. \textbf{Across the three judge pairs, Cohen's κ
falls in {[}−0.012, 0.000{]} (two pairs degenerate; the third less than
chance) while Gwet's AC1 lands in {[}0.961, 0.992{]} --- almost perfect
agreement on all three pairs.} Appendix A.3 lists the full per-pair κ,
AC1, and Landis \& Koch interpretation.

\textbf{Five disagreement cases, two distinct patterns.} The 2-judge
headline saw a single disagreement; the 3-judge headline sees five.
Decomposing:

\begin{longtable}[]{@{}
  >{\raggedright\arraybackslash}p{(\linewidth - 12\tabcolsep) * \real{0.1429}}
  >{\raggedright\arraybackslash}p{(\linewidth - 12\tabcolsep) * \real{0.1429}}
  >{\raggedright\arraybackslash}p{(\linewidth - 12\tabcolsep) * \real{0.1429}}
  >{\raggedright\arraybackslash}p{(\linewidth - 12\tabcolsep) * \real{0.1429}}
  >{\raggedright\arraybackslash}p{(\linewidth - 12\tabcolsep) * \real{0.1429}}
  >{\raggedright\arraybackslash}p{(\linewidth - 12\tabcolsep) * \real{0.1429}}
  >{\raggedright\arraybackslash}p{(\linewidth - 12\tabcolsep) * \real{0.1429}}@{}}
\toprule\noalign{}
\begin{minipage}[b]{\linewidth}\raggedright
Case
\end{minipage} & \begin{minipage}[b]{\linewidth}\raggedright
Tier
\end{minipage} & \begin{minipage}[b]{\linewidth}\raggedright
Domain
\end{minipage} & \begin{minipage}[b]{\linewidth}\raggedright
Primary
\end{minipage} & \begin{minipage}[b]{\linewidth}\raggedright
Secondary
\end{minipage} & \begin{minipage}[b]{\linewidth}\raggedright
Tertiary
\end{minipage} & \begin{minipage}[b]{\linewidth}\raggedright
Pattern
\end{minipage} \\
\midrule\noalign{}
\endhead
\bottomrule\noalign{}
\endlastfoot
\texttt{cardio\_17} & 1 & cardio & FAITHFUL & FAITHFUL & \textbf{HALL} &
tertiary-only \\
\texttt{cardio\_40} & 2 & cardio & FAITHFUL & \textbf{HALL} & FAITHFUL &
secondary-only (historical case from the prior 2-specialty
configuration; §6.3) \\
\texttt{endo\_21} & 1 & endo & FAITHFUL & FAITHFUL & \textbf{HALL} &
tertiary-only \\
\texttt{endo\_38} & 2 & endo & FAITHFUL & FAITHFUL & \textbf{HALL} &
tertiary-only \\
\texttt{infect\_34} & 1 & infect & FAITHFUL & FAITHFUL & \textbf{HALL} &
tertiary-only \\
\end{longtable}

The four tertiary-only flags are the empirical signature of same-family
bias: a model from a different vendor (OpenAI rather than Yandex) said
HALLUCINATION on cases the two Yandex judges both said FAITHFUL. Under
the minimum-judge rule the test-split tier breakdown is \textbf{100.0\%
on every (domain, tier) cell except T1 cardiology, T2 cardiology, T1
endocrinology, T2 endocrinology, and T1 infectiology, where one case
each disagrees}. Per-domain minimum-judge rates: cardiology 33/35 =
94.3\% {[}81.4\%--98.4\%{]}, endocrinology 33/35 = 94.3\%
{[}81.4\%--98.4\%{]}, gastroenterology 32/32 = 100.0\%
{[}89.3\%--100.0\%{]}, infectiology 29/30 = 96.7\% {[}83.3\%--99.4\%{]}.

\textbf{Regression check.} The single Yandex-pair disagreement is
\texttt{cardio\_40} in both the n = 70 prior 2-specialty run and the n =
140 4-specialty run --- the new gastro/infect cases produced no
additional disagreements between the two Yandex judges. The
2-Yandex-judge minimum rate rose from 69/70 = 98.6\% (prior 2-specialty)
to 131/132 = 99.2\% (4-specialty) simply because the denominator nearly
doubled while the Yandex disagreement count stayed at 1. The 3-judge
minimum-judge rate of \textbf{127/132 = 96.2\%} is \textbf{−3.0 pp below
the 2-Yandex-judge upper bound} of 99.2\%; the gap is the cross-vendor
judge's contribution. The observed −3.0 pp drop is at the low end of the
5--25 pp same-family inflation range Zheng et al.~characterise
(\citeproc{ref-zheng2023mtbench}{Zheng et al., 2023}); the high baseline
(99.2\%) leaves limited room for upward inflation, which is consistent
with a smaller observed effect than the midpoint of the predicted range.
The 2-judge upper bound remains valid as the \emph{upper} faithfulness
bound under same-family bias; the 3-judge result is now the
production-quality \emph{lower} bound, with a Wilson lower endpoint of
91.4\%.

Run cost reference. \textbf{2-judge stage (2026-05-21):} 266 Yandex
calls (140 × 2 minus 14 fallback-skipped), 15.2 min wall-clock. Raw
per-case verdicts from the 2-judge stage are retained so the
tertiary-only driver can reuse the Yandex pair without re-paying.
\textbf{3-judge stage (2026-05-23):} 133 OpenRouter GPT-OSS-120B calls
(free tier), 29.4 min wall-clock, 0 errors. An earlier multijudge
reconciliation note documents the variance pattern at temperature=0 on
the smaller yandexgpt-lite model; a 4-specialty reconciliation update is
a follow-up flagged in §7 L8.

\subsection{5.5 Out-of-Scope Refusal
Gate}\label{out-of-scope-refusal-gate}

Architectural framing. FAISS always returns K nearest neighbours by
construction --- it cannot refuse. Two L2 distance thresholds operate on
the same scalar in opposite directions and with conflicting objectives:
\texttt{MAX\_L2\_DISTANCE\ =\ 1.2} is the retrieval quality filter,
tuned on the dev split to maximise in-corpus Recall@5;
\texttt{L2\_REJECT\_MIN} is the refusal gate, tuned on the same min-L2
distribution to maximise out-of-scope refusal. Because the in-corpus
T1/T2 min-L2 distribution overlaps the out-of-scope T3 distribution on
this corpus and embedding model, no single scalar threshold separates
them; any choice trades T3 recall against T1/T2 false-positive rate.

\textbf{Re-tune for the 4-specialty corpus.} The prior 2-specialty
configuration used \texttt{L2\_REJECT\_MIN\ =\ 0.92}. On the new
4-specialty data, the in-corpus mean L2 (μ) is 0.874 (cardio), 0.894
(endo), 0.880 (gastro), 0.894 (infect) --- i.e.~the new corpora's μ is
\textbf{above} 0.92 for endo and infect, so the old threshold catches an
unacceptable 71.2\% (79/111) of in-corpus T1/T2 queries as false
positives. The refusal-gate tuner was re-run on the new 60-case dev
split (Signal A grid + Signal B grid). No threshold in the grid
simultaneously hits the prior target of
\texttt{≥80\%\ T3\ recall\ AND\ ≤5\%\ T1/T2\ FP\ rate} on dev --- the
in-corpus / out-of-corpus min-L2 distributions overlap too heavily. The
tuner falls back to the best-F1 row, yielding
\textbf{\texttt{L2\_REJECT\_MIN\ =\ 1.020}}. The trade-off curve has
shifted dramatically:

\textbf{Chosen signal: A (minimum-L2 threshold). Chosen threshold:
\texttt{L2\_REJECT\_MIN\ =\ 1.020}.} The refusal-gate module is invoked
from the specialist module \emph{before} the LLM call (the generic
single-class implementation is unchanged in this evaluation). If
\texttt{min(L2\ distances\ over\ top-K=5\ retrieved\ chunks)\ \textgreater{}\ L2\_REJECT\_MIN},
the agent short-circuits and returns the canned ``Insufficient evidence
in the current knowledge base to address this specific query.'' response
without calling the generation model.

\subsubsection{Per-specialty in-corpus / out-of-corpus L2 ranges
(140-case test
split)}\label{per-specialty-in-corpus-out-of-corpus-l2-ranges-140-case-test-split}

\begin{longtable}[]{@{}
  >{\raggedright\arraybackslash}p{(\linewidth - 6\tabcolsep) * \real{0.2500}}
  >{\raggedright\arraybackslash}p{(\linewidth - 6\tabcolsep) * \real{0.2500}}
  >{\raggedright\arraybackslash}p{(\linewidth - 6\tabcolsep) * \real{0.2500}}
  >{\raggedright\arraybackslash}p{(\linewidth - 6\tabcolsep) * \real{0.2500}}@{}}
\toprule\noalign{}
\begin{minipage}[b]{\linewidth}\raggedright
Specialty
\end{minipage} & \begin{minipage}[b]{\linewidth}\raggedright
T1/T2 min-L2 (n)
\end{minipage} & \begin{minipage}[b]{\linewidth}\raggedright
T3 min-L2 (n)
\end{minipage} & \begin{minipage}[b]{\linewidth}\raggedright
Overlap zone
\end{minipage} \\
\midrule\noalign{}
\endhead
\bottomrule\noalign{}
\endlastfoot
Cardiology & {[}0.831, 1.073{]} (n = 27) & {[}0.906, 1.000{]} (n = 8) &
{[}0.906, 1.000{]} \\
Endocrinology & {[}0.797, 1.008{]} (n = 28) & {[}0.844, 0.973{]} (n = 7)
& {[}0.844, 0.973{]} \\
Gastroenterology & {[}0.874, 1.071{]} (n = 28) & {[}0.970, 1.063{]} (n =
7) & {[}0.970, 1.063{]} \\
Infectiology & {[}0.963, 1.107{]} (n = 28) & {[}0.977, 1.117{]} (n = 7)
& {[}0.977, 1.107{]} \\
\textbf{All 4 specialties} & \textbf{{[}0.797, 1.107{]} (n = 111)} &
\textbf{{[}0.844, 1.117{]} (n = 29)} & \textbf{{[}0.844, 1.107{]}} \\
\end{longtable}

Infectiology shows nearly complete overlap (T1/T2 range and T3 range
have only a 0.014 separation at the top). Cardiology has the cleanest
separation. An earlier architectural framing (``T1/T2 0.70--1.07; T3
0.84--1.00'') was for the 2-specialty corpus; the 4-specialty ranges
above replace it.

\subsubsection{\texorpdfstring{Test-split precision / recall at
\texttt{L2\_REJECT\_MIN\ =\ 1.020}}{Test-split precision / recall at L2\_REJECT\_MIN = 1.020}}\label{test-split-precision-recall-at-l2_reject_min-1.020}

Positive class = Tier 3 (correct outcome: refuse). Negative class = Tier
1/2 (correct outcome: pass through to the LLM).

\begin{longtable}[]{@{}
  >{\raggedright\arraybackslash}p{(\linewidth - 8\tabcolsep) * \real{0.2000}}
  >{\raggedright\arraybackslash}p{(\linewidth - 8\tabcolsep) * \real{0.2000}}
  >{\raggedright\arraybackslash}p{(\linewidth - 8\tabcolsep) * \real{0.2000}}
  >{\raggedright\arraybackslash}p{(\linewidth - 8\tabcolsep) * \real{0.2000}}
  >{\raggedright\arraybackslash}p{(\linewidth - 8\tabcolsep) * \real{0.2000}}@{}}
\toprule\noalign{}
\begin{minipage}[b]{\linewidth}\raggedright
Stratum
\end{minipage} & \begin{minipage}[b]{\linewidth}\raggedright
Cases
\end{minipage} & \begin{minipage}[b]{\linewidth}\raggedright
Refused by gate
\end{minipage} & \begin{minipage}[b]{\linewidth}\raggedright
Refusal rate
\end{minipage} & \begin{minipage}[b]{\linewidth}\raggedright
Wilson 95\% CI
\end{minipage} \\
\midrule\noalign{}
\endhead
\bottomrule\noalign{}
\endlastfoot
Tier 3 (positive class) & 29 & \textbf{6} & \textbf{20.7\%} &
{[}9.8\%--38.4\%{]} \\
Tier 1/2 (negative class --- FP) & 111 & \textbf{19} & \textbf{17.1\%}
(FP rate) & {[}11.2\%--25.2\%{]} \\
\end{longtable}

Appendix A.4 contains the full per-Tier per-specialty refusal breakdown
(12 cells across cardiology / endocrinology / gastroenterology /
infectiology × T1 / T2 / T3).

\subsubsection{Target check}\label{target-check}

\begin{longtable}[]{@{}lll@{}}
\toprule\noalign{}
Target & Achieved? & Numbers \\
\midrule\noalign{}
\endhead
\bottomrule\noalign{}
\endlastfoot
≥80\% Tier 3 rejection on test & \textbf{✗} & 6/29 = 20.7\% \\
≤5\% Tier 1/2 FP rate on test & \textbf{✗} & 19/111 = 17.1\% \\
\end{longtable}

Both targets are missed at the chosen threshold. The original target
ratios (80\% recall, 5\% FP) \textbf{cannot be simultaneously met by any
single-threshold gate on the 4-specialty corpus} --- the best F1 falls
well below either target.

\subsubsection{What the trade-off looks like on the 4-specialty
corpus}\label{what-the-trade-off-looks-like-on-the-4-specialty-corpus}

\begin{longtable}[]{@{}
  >{\raggedright\arraybackslash}p{(\linewidth - 6\tabcolsep) * \real{0.2500}}
  >{\raggedright\arraybackslash}p{(\linewidth - 6\tabcolsep) * \real{0.2500}}
  >{\raggedright\arraybackslash}p{(\linewidth - 6\tabcolsep) * \real{0.2500}}
  >{\raggedright\arraybackslash}p{(\linewidth - 6\tabcolsep) * \real{0.2500}}@{}}
\toprule\noalign{}
\begin{minipage}[b]{\linewidth}\raggedright
\texttt{L2\_REJECT\_MIN}
\end{minipage} & \begin{minipage}[b]{\linewidth}\raggedright
Test T3 recall
\end{minipage} & \begin{minipage}[b]{\linewidth}\raggedright
Test T1/T2 FP rate
\end{minipage} & \begin{minipage}[b]{\linewidth}\raggedright
Comment
\end{minipage} \\
\midrule\noalign{}
\endhead
\bottomrule\noalign{}
\endlastfoot
0.900 & 96.6\% (28/29) & 78.4\% (87/111) & Prior gate minimum;
over-refuses Tier 1/2 \\
\textbf{0.920 (prior gate)} & \textbf{89.7\% (26/29)} & \textbf{71.2\%
(79/111)} & Old threshold, broken on 4 specialties \\
\textbf{1.020 (current)} & \textbf{20.7\% (6/29)} & \textbf{17.1\%
(19/111)} & Best-F1 fallback at current rule \\
1.100 & 3.4\% (1/29) & 2.7\% (3/111) & Near-zero refusal \\
1.200 (= MAX\_L2) & 0.0\% (0/29) & 0.0\% (0/111) & Gate effectively
off \\
\end{longtable}

\begin{figure}
\centering
\pandocbounded{\includesvg[keepaspectratio]{figures/F2_refusal_sweep.svg}}
\caption{Tier 3 refusal recall vs Tier 1/2 false-positive rate as the
numeric refusal gate L2\_REJECT\_MIN sweeps across the tested range on
the held-out test split. The chosen threshold (1.020) does not
simultaneously meet both target rates.}
\end{figure}

\subsubsection{Comparison to the prior prompt-only refusal and the prior
numeric
gate}\label{comparison-to-the-prior-prompt-only-refusal-and-the-prior-numeric-gate}

\begin{longtable}[]{@{}
  >{\raggedright\arraybackslash}p{(\linewidth - 6\tabcolsep) * \real{0.2500}}
  >{\raggedright\arraybackslash}p{(\linewidth - 6\tabcolsep) * \real{0.2500}}
  >{\raggedright\arraybackslash}p{(\linewidth - 6\tabcolsep) * \real{0.2500}}
  >{\raggedright\arraybackslash}p{(\linewidth - 6\tabcolsep) * \real{0.2500}}@{}}
\toprule\noalign{}
\begin{minipage}[b]{\linewidth}\raggedright
Metric
\end{minipage} & \begin{minipage}[b]{\linewidth}\raggedright
Prompt-only (§6.2 baseline)
\end{minipage} & \begin{minipage}[b]{\linewidth}\raggedright
Prior numeric gate (\texttt{L2\_REJECT\_MIN=0.92})
\end{minipage} & \begin{minipage}[b]{\linewidth}\raggedright
Current numeric gate (\texttt{L2\_REJECT\_MIN=1.020})
\end{minipage} \\
\midrule\noalign{}
\endhead
\bottomrule\noalign{}
\endlastfoot
Test Tier 3 rejection rate & 0/29 (0.0\%) & 26/29 (89.7\%) & 6/29
(20.7\%) \\
Test Tier 1/2 FP rate & 0/111 (0.0\%) & 79/111 (71.2\%) & 19/111
(17.1\%) \\
\end{longtable}

The 0.92 → 1.020 re-tune trades \textbf{a 69 pp loss in T3 recall} for
\textbf{a 54 pp drop in T1/T2 FP rate}. Both are deliberate but neither
is good: at 0.92 the gate over-refused two-thirds of legitimate queries;
at 1.020 it lets through nearly all of them but also catches only a
fifth of true T3 cases. \textbf{The two-stage gate (L2 pre-filter
feeding an LLM-as-classifier confirmer), previously flagged as future
work, is now urgent}, not optional. See §7 Limitation 7 for the proposed
design.

\subsection{5.6 Offline Retrieval Regression
Test}\label{offline-retrieval-regression-test}

To guard against silent retrieval drift (threshold changes, index
corruption, accidental re-embedding) without burning Yandex API calls on
every CI run, an offline regression test is included. \textbf{Twenty
representative queries} (5 per specialty: cardiology, endocrinology,
gastroenterology, infectiology) are pre-embedded once via the live
Yandex API and saved as a snapshot. Subsequent test runs load the saved
vectors and call FAISS search directly on the binary indices, bypassing
both LangChain and the embedding model API. The test asserts (a) every
query's top-K=5 source-file set equals the snapshot's set; (b) per-rank
L2 drift is \textless{} 0.1; (c) every query retrieves at least one
chunk within \texttt{MAX\_L2\_DISTANCE\ =\ 1.2}. An earlier extension
from 10 → 20 queries doubled the regression coverage; the cardio/endo
subset is byte-identical to the original 2-specialty baseline as a
verified regression check.

This is a regression check, not a new evaluation metric --- it does not
affect the numbers reported in §5.1--§5.5.

\subsection{5.7 Summary of All Metrics (200-Case Tiered Dataset, 4
specialties)}\label{summary-of-all-metrics-200-case-tiered-dataset-4-specialties}

The metrics below are broken down by domain and difficulty tier on the
\textbf{full 200-case dataset} (4 specialties × 50 cases). Numbers in
this table aggregate across dev (60 cases, used for hyperparameter
tuning) and test (140 cases, held-out). The held-out test-split-only
table is in §5.8. The Retrieval row reports \textbf{Recall@5} (primary
grounded metric) with the legacy KeywordHitRate next to it; Tier 3
measures safety-fallback behaviour rather than retrieval hit rate.

\begin{longtable}[]{@{}
  >{\raggedright\arraybackslash}p{(\linewidth - 18\tabcolsep) * \real{0.1000}}
  >{\raggedright\arraybackslash}p{(\linewidth - 18\tabcolsep) * \real{0.1000}}
  >{\raggedright\arraybackslash}p{(\linewidth - 18\tabcolsep) * \real{0.1000}}
  >{\raggedright\arraybackslash}p{(\linewidth - 18\tabcolsep) * \real{0.1000}}
  >{\raggedright\arraybackslash}p{(\linewidth - 18\tabcolsep) * \real{0.1000}}
  >{\raggedright\arraybackslash}p{(\linewidth - 18\tabcolsep) * \real{0.1000}}
  >{\raggedright\arraybackslash}p{(\linewidth - 18\tabcolsep) * \real{0.1000}}
  >{\raggedright\arraybackslash}p{(\linewidth - 18\tabcolsep) * \real{0.1000}}
  >{\raggedright\arraybackslash}p{(\linewidth - 18\tabcolsep) * \real{0.1000}}
  >{\raggedright\arraybackslash}p{(\linewidth - 18\tabcolsep) * \real{0.1000}}@{}}
\toprule\noalign{}
\begin{minipage}[b]{\linewidth}\raggedright
Metric
\end{minipage} & \begin{minipage}[b]{\linewidth}\raggedright
T1 Cardio
\end{minipage} & \begin{minipage}[b]{\linewidth}\raggedright
T2 Cardio
\end{minipage} & \begin{minipage}[b]{\linewidth}\raggedright
T1 Endo
\end{minipage} & \begin{minipage}[b]{\linewidth}\raggedright
T2 Endo
\end{minipage} & \begin{minipage}[b]{\linewidth}\raggedright
T1 Gastro
\end{minipage} & \begin{minipage}[b]{\linewidth}\raggedright
T2 Gastro
\end{minipage} & \begin{minipage}[b]{\linewidth}\raggedright
T1 Infect
\end{minipage} & \begin{minipage}[b]{\linewidth}\raggedright
T2 Infect
\end{minipage} & \begin{minipage}[b]{\linewidth}\raggedright
T3 (all 4)
\end{minipage} \\
\midrule\noalign{}
\endhead
\bottomrule\noalign{}
\endlastfoot
Routing Accuracy & 100.0\% {[}87.0--100\%{]} & 100.0\% {[}78.5--100\%{]}
& 100.0\% {[}87.0--100\%{]} & 100.0\% {[}80.6--100\%{]} & 88.9\% (24/27)
{[}71.9--96.2\%{]} & 86.7\% (13/15) {[}62.1--96.3\%{]} & 96.3\% (26/27)
{[}81.7--99.3\%{]} & 100.0\% {[}79.6--100\%{]} & 100.0\% (32/32)
{[}89.3--100\%{]} \\
Retrieval Recall@5 (gold-doc Bernoulli) & 64.2\% (52/81)
{[}53.3--73.8\%{]} & 54.1\% (20/37) {[}38.4--69.0\%{]} & 60.3\% (47/78)
{[}49.2--70.4\%{]} & 52.3\% (23/44) {[}37.9--66.2\%{]} & 57.5\% (42/73)
{[}46.1--68.2\%{]} & 60.0\% (21/35) {[}43.6--74.4\%{]} & 51.5\% (35/68)
{[}39.8--62.9\%{]} & 62.5\% (20/32) {[}45.3--77.1\%{]} & \emph{Refusal
rate --- see §5.3} \\
Retrieval KeywordHitRate (legacy) & 100.0\% {[}87.0--100\%{]} & 78.6\%
(11/14) {[}52.4--92.4\%{]} & 96.3\% (26/27) {[}81.7--99.3\%{]} & 93.8\%
(15/16) {[}71.7--98.9\%{]} & 81.5\% (22/27) {[}63.3--92.1\%{]} & 80.0\%
(12/15) {[}54.8--93.0\%{]} & 85.2\% (23/27) {[}67.5--94.1\%{]} & 73.3\%
(11/15) {[}48.0--89.1\%{]} & \emph{See §5.3 note on adjacent content} \\
Faithfulness (minimum-judge, n = 140 test only --- §5.4) & 100.0\%
(13/13) {[}77.2--100\%{]} & 92.9\% (13/14) {[}68.5--98.7\%{]} & 100.0\%
(12/12) {[}75.8--100\%{]} & 100.0\% (16/16) {[}80.6--100\%{]} & 100.0\%
(13/13) {[}77.2--100\%{]} & 100.0\% (15/15) {[}79.6--100\%{]} & 100.0\%
(12/12) {[}75.8--100\%{]} & 100.0\% (15/15) {[}79.6--100\%{]} & 100.0\%
(22/22) {[}85.1--100\%{]} \\
\end{longtable}

\emph{(Routing / Retrieval / KeywordHitRate cells are computed on the
full 200-case dataset. Faithfulness is the minimum-judge rate on the
140-case test split only --- the cardio\_40 disagreement (the single
primary/secondary disagreement across all 132 judged cases) sits in T2
cardiology and drops that cell to 13/14 = 92.9\%; every other cell is
100\% by minimum-judge. Tier 3 faithfulness counts only T3 cases that
did NOT trigger the refusal-gate fallback (22 of the 29 test-split T3
cases). Confidence intervals are 95\% Wilson via \texttt{statsmodels}.
Recall@5 denominators are gold-doc-level Bernoulli trials.)}

Under the legacy KeywordHitRate, T1 cardiology reads 100\% --- but the
grounded Recall@5 on the same cases is 64.2\%, so retrieval surfaces
only \textasciitilde2 of the 3 gold documents in the top-5 window for
the average T1 cardiology query. The 4-specialty extension preserves
this Recall@5-vs-KeywordHitRate gap on every domain. Routing stays at
≥86.7\% across every tier × domain; faithfulness stays at 100\% by
minimum-judge on every cell except T2 cardiology (cardio\_40
disagreement). The refusal gate (§5.5) catches 20.7\% (6/29) of Tier-3
test cases at the cost of a 17.1\% Tier-1/2 false-positive rate --- the
trade-off has worsened vs the prior single-threshold-gate baseline on
4-specialty data and is the subject of an open work item (§7 L7).

\subsection{5.8 Held-Out Test Set Results (n =
140)}\label{held-out-test-set-results-n-140}

To provide an unbiased measurement of generalisation, all evaluations
were re-run on the 140-case held-out test split, which contains every
case in the golden dataset except the 60 development cases (cases 1--15
from each of the four specialties) used for hyperparameter tuning. Tier
composition of the test split: T1 = 51 (13 cardio + 12 endo + 13 gastro
+ 13 infect), T2 = 60 (14 cardio + 16 endo + 15 gastro + 15 infect), T3
= 29 (8 cardio + 7 endo + 7 gastro + 7 infect).

\subsubsection{Retrieval Hit Rate (Test Split, KeywordHitRate legacy +
Recall@5
grounded)}\label{retrieval-hit-rate-test-split-keywordhitrate-legacy-recall5-grounded}

\begin{longtable}[]{@{}
  >{\raggedright\arraybackslash}p{(\linewidth - 6\tabcolsep) * \real{0.2500}}
  >{\raggedright\arraybackslash}p{(\linewidth - 6\tabcolsep) * \real{0.2500}}
  >{\raggedright\arraybackslash}p{(\linewidth - 6\tabcolsep) * \real{0.2500}}
  >{\raggedright\arraybackslash}p{(\linewidth - 6\tabcolsep) * \real{0.2500}}@{}}
\toprule\noalign{}
\begin{minipage}[b]{\linewidth}\raggedright
Domain
\end{minipage} & \begin{minipage}[b]{\linewidth}\raggedright
KeywordHitRate (legacy)
\end{minipage} & \begin{minipage}[b]{\linewidth}\raggedright
Recall@5 (grounded)
\end{minipage} & \begin{minipage}[b]{\linewidth}\raggedright
Random Hit Rate
\end{minipage} \\
\midrule\noalign{}
\endhead
\bottomrule\noalign{}
\endlastfoot
Cardiology & 88.6\% (31/35) {[}74.0\%--95.5\%{]} & 56.6\% (43/76)
{[}45.4\%--67.1\%{]} & 25.7\% (9/35) {[}14.2\%--42.1\%{]} \\
Endocrinology & 94.3\% (33/35) {[}81.4\%--98.4\%{]} & 55.8\% (43/77)
{[}44.6\%--66.5\%{]} & 20.0\% (7/35) {[}10.0\%--35.9\%{]} \\
Gastroenterology & 80.0\% (28/35) {[}64.1\%--90.0\%{]} & 56.3\% (40/71)
{[}44.8\%--67.3\%{]} & 5.7\% (2/35) {[}1.6\%--18.6\%{]} \\
Infectiology & 80.0\% (28/35) {[}64.1\%--90.0\%{]} & 56.3\% (36/64)
{[}44.1\%--67.8\%{]} & 0.0\% (0/35) {[}0.0\%--9.9\%{]} \\
\textbf{Overall} & \textbf{85.7\% (120/140) {[}78.9\%--90.6\%{]}} &
\textbf{56.2\% (162/288) {[}50.5\%--61.9\%{]}} & \textbf{12.9\% (18/140)
{[}8.3\%--19.5\%{]}} \\
\end{longtable}

\subsubsection{Summary of All Metrics (Test Split, n =
140)}\label{summary-of-all-metrics-test-split-n-140}

\begin{longtable}[]{@{}
  >{\raggedright\arraybackslash}p{(\linewidth - 18\tabcolsep) * \real{0.1000}}
  >{\raggedright\arraybackslash}p{(\linewidth - 18\tabcolsep) * \real{0.1000}}
  >{\raggedright\arraybackslash}p{(\linewidth - 18\tabcolsep) * \real{0.1000}}
  >{\raggedright\arraybackslash}p{(\linewidth - 18\tabcolsep) * \real{0.1000}}
  >{\raggedright\arraybackslash}p{(\linewidth - 18\tabcolsep) * \real{0.1000}}
  >{\raggedright\arraybackslash}p{(\linewidth - 18\tabcolsep) * \real{0.1000}}
  >{\raggedright\arraybackslash}p{(\linewidth - 18\tabcolsep) * \real{0.1000}}
  >{\raggedright\arraybackslash}p{(\linewidth - 18\tabcolsep) * \real{0.1000}}
  >{\raggedright\arraybackslash}p{(\linewidth - 18\tabcolsep) * \real{0.1000}}
  >{\raggedright\arraybackslash}p{(\linewidth - 18\tabcolsep) * \real{0.1000}}@{}}
\toprule\noalign{}
\begin{minipage}[b]{\linewidth}\raggedright
Metric
\end{minipage} & \begin{minipage}[b]{\linewidth}\raggedright
T1 Cardio
\end{minipage} & \begin{minipage}[b]{\linewidth}\raggedright
T2 Cardio
\end{minipage} & \begin{minipage}[b]{\linewidth}\raggedright
T1 Endo
\end{minipage} & \begin{minipage}[b]{\linewidth}\raggedright
T2 Endo
\end{minipage} & \begin{minipage}[b]{\linewidth}\raggedright
T1 Gastro
\end{minipage} & \begin{minipage}[b]{\linewidth}\raggedright
T2 Gastro
\end{minipage} & \begin{minipage}[b]{\linewidth}\raggedright
T1 Infect
\end{minipage} & \begin{minipage}[b]{\linewidth}\raggedright
T2 Infect
\end{minipage} & \begin{minipage}[b]{\linewidth}\raggedright
T3 (all 4)
\end{minipage} \\
\midrule\noalign{}
\endhead
\bottomrule\noalign{}
\endlastfoot
Routing Accuracy & 100.0\% (13/13) {[}77.2--100\%{]} & 100.0\% (14/14)
{[}78.5--100\%{]} & 100.0\% (12/12) {[}75.8--100\%{]} & 93.8\% (15/16)
{[}71.7--98.9\%{]} & 84.6\% (11/13) {[}57.8--95.7\%{]} & 86.7\% (13/15)
{[}62.1--96.3\%{]} & 92.3\% (12/13) {[}66.7--98.6\%{]} & 100.0\% (15/15)
{[}79.6--100\%{]} & 100.0\% (29/29) {[}88.3--100\%{]} \\
Retrieval Recall@5 (pooled gold-doc) & 59.0\% (23/39) {[}43.4--72.9\%{]}
& 54.1\% (20/37) {[}38.4--69.0\%{]} & 60.6\% (20/33) {[}43.7--75.3\%{]}
& 52.3\% (23/44) {[}37.9--66.2\%{]} & 52.8\% (19/36) {[}37.0--68.0\%{]}
& 60.0\% (21/35) {[}43.6--74.4\%{]} & 50.0\% (16/32) {[}33.6--66.4\%{]}
& 62.5\% (20/32) {[}45.3--77.1\%{]} & \emph{Refusal Rate --- see
§5.5} \\
Retrieval KeywordHitRate (legacy) & 100.0\% (13/13) {[}77.2--100\%{]} &
78.6\% (11/14) {[}52.4--92.4\%{]} & 91.7\% (11/12) {[}64.6--98.5\%{]} &
93.8\% (15/16) {[}71.7--98.9\%{]} & 76.9\% (10/13) {[}49.7--91.8\%{]} &
80.0\% (12/15) {[}54.8--93.0\%{]} & 84.6\% (11/13) {[}57.8--95.7\%{]} &
73.3\% (11/15) {[}48.0--89.1\%{]} & \emph{See §5.3 note} \\
Faithfulness (minimum-judge) & 100.0\% (13/13) {[}77.2--100\%{]} &
92.9\% (13/14) {[}68.5--98.7\%{]} & 100.0\% (12/12) {[}75.8--100\%{]} &
100.0\% (16/16) {[}80.6--100\%{]} & 100.0\% (13/13) {[}77.2--100\%{]} &
100.0\% (15/15) {[}79.6--100\%{]} & 100.0\% (12/12) {[}75.8--100\%{]} &
100.0\% (15/15) {[}79.6--100\%{]} & 100.0\% (22/22) {[}85.1--100\%{]} \\
Chunk-relevance (LLM judge) & 100.0\% {[}77.2--100\%{]} & 100.0\%
{[}78.5--100\%{]} & 100.0\% {[}75.8--100\%{]} & 93.8\% (15/16)
{[}71.7--98.9\%{]} & 100.0\% {[}77.2--100\%{]} & 100.0\%
{[}79.6--100\%{]} & 100.0\% {[}75.8--100\%{]} & 100.0\%
{[}79.6--100\%{]} & 100.0\% (29/29) {[}88.3--100\%{]} \\
\end{longtable}

\emph{(Wilson 95\% CIs via \texttt{statsmodels}. The single
chunk-relevance miss is in T2 endocrinology --- same case
(\texttt{endo\_25}) that has surfaced in prior stages as a borderline
case. The single faithfulness miss is \texttt{cardio\_40} (T2 cardio ---
congenital LQTS). Of the 8 cases excluded from the 132
judge-intersection denominator: 6 T3 cases triggered the refusal gate (3
gastro + 3 infect; no answer generated, no faithfulness assessment
possible), 1 further infect T3 case (\texttt{infect\_17}) returned a
None verdict due to a primary-judge API failure, and 1 case was excluded
by the prompt-only Insufficient evidence fallback. Total excluded: 6
gate fires + 1 API failure + 1 prompt fallback = 8; judged denominator
140 − 8 = 132.)}

\textbf{Headline:} Routing accuracy on test = \textbf{95.7\% (134/140)
{[}91.0\%--98.0\%{]}}; FAISS Recall@5 = \textbf{56.2\% (162/288)
{[}50.5\%--61.9\%{]}}; Minimum-judge faithfulness = \textbf{96.2\%
(127/132) {[}91.4\%--98.4\%{]}} (3-judge, cross-vendor lower bound;
2-Yandex-judge upper bound = 99.2\% {[}95.8\%--99.9\%{]}); Numeric-gate
T3 refusal at the new threshold = \textbf{20.7\% (6/29)
{[}9.8\%--38.4\%{]}} with T1/T2 FP rate = 17.1\%.

Across all four specialties the held-out split confirms the refusal gate
as the sole unresolved failure mode; every other metric lands within a
few percentage points of its full-set value.

\subsection{5.9 External Benchmark: PubMedQA Cardiology
Slice}\label{external-benchmark-pubmedqa-cardiology-slice}

To anchor the in-house Recall@5 against an independently labelled
biomedical retrieval benchmark --- addressing the auto-annotation
circularity disclosed in §5.3 --- the cardiologist agent's FAISS index
is evaluated against PubMedQA's expert-labelled subset
(\citeproc{ref-jin2019pubmedqa}{Jin et al., 2019}), downloadable from
HuggingFace as \texttt{qiaojin/PubMedQA}, subset \texttt{pqa\_labeled}
(1000 manually curated yes/no/maybe research-question QA pairs).
Filtering the 1000-case split to cardiology-relevant questions via a
case-insensitive substring OR over \{\texttt{heart}, \texttt{cardiac},
\texttt{cardio}, \texttt{ventricular}, \texttt{atrial},
\texttt{coronary}, \texttt{mitral}, \texttt{aortic}, \texttt{valve},
\texttt{arrhythmia}, \texttt{hypertension}, \texttt{stroke}\} yields
\textbf{n = 85 questions} with 275 gold abstract passages across them.

\begin{longtable}[]{@{}
  >{\raggedright\arraybackslash}p{(\linewidth - 6\tabcolsep) * \real{0.2500}}
  >{\raggedright\arraybackslash}p{(\linewidth - 6\tabcolsep) * \real{0.2500}}
  >{\raggedright\arraybackslash}p{(\linewidth - 6\tabcolsep) * \real{0.2500}}
  >{\raggedright\arraybackslash}p{(\linewidth - 6\tabcolsep) * \real{0.2500}}@{}}
\toprule\noalign{}
\begin{minipage}[b]{\linewidth}\raggedright
Source
\end{minipage} & \begin{minipage}[b]{\linewidth}\raggedright
Recall@5 (pooled)
\end{minipage} & \begin{minipage}[b]{\linewidth}\raggedright
n (gold trials)
\end{minipage} & \begin{minipage}[b]{\linewidth}\raggedright
95\% Wilson CI
\end{minipage} \\
\midrule\noalign{}
\endhead
\bottomrule\noalign{}
\endlastfoot
This work (in-house, held-out test split, cardiology) & 56.6\% (43/76) &
76 gold-doc Bernoulli trials & {[}45.4\%--67.1\%{]} \\
PubMedQA cardiology slice (sentence-level Jaccard ≥ 0.20) & 21.5\%
(59/275) & 275 gold-passage Bernoulli trials & {[}17.0\%--26.7\%{]} \\
\end{longtable}

Matching threshold: each retrieved chunk and each gold passage is split
into sentences on \texttt{{[}.!?{]}} boundaries, tokens are lowercased
alphanumeric words of length ≥ 2, and a chunk is judged to \emph{hit} a
gold passage when at least one (chunk\_sentence, gold\_sentence) pair
reaches token-level Jaccard
\texttt{\textbar{}A\ ∩\ B\textbar{}\ /\ \textbar{}A\ ∪\ B\textbar{}} ≥
0.20. The spec's preferred threshold (≥ 0.30) was empirically
unreachable on this corpus pair --- a probe across all 275 gold passages
found the maximum achievable sentence-pair Jaccard was 0.294 (mean
0.163), because the cardiology corpus is written in clinical-guideline /
textbook register while PubMedQA passages are research-abstract
register. 0.20 sits at the 21.5\% percentile of the achievable
distribution and is the operating point that surfaces a non-zero
comparison signal without being dominated by stopword overlap. PubMedQA
itself (Jin et al.~2019 (\citeproc{ref-jin2019pubmedqa}{Jin et al.,
2019})) does not define a canonical Jaccard threshold for retrieval
matching --- it uses BERT-based reading-comprehension evaluation against
a single labelled answer. The Jaccard-based matching here is a
deliberately simple lexical surrogate chosen so the per-passage hit rule
is reproducible without any judge LLM, accepting that it under-counts
semantically correct retrievals that paraphrase rather than lexically
overlap. The two rows in the table above are therefore not directly
comparable: the in-house row uses doc-level identity matching against
gold sources auto-annotated from the same retrieval system (the very
circularity disclosed in §5.3), whereas the PubMedQA row uses lexical
Jaccard matching against an independently labelled corpus from a
different register entirely. The 35-point gap is consistent with both
interpretations --- (a) the in-house number is inflated by the
same-FAISS-system gold-source bias, and (b) the PubMedQA Jaccard rule
under-counts paraphrastic matches --- and the two contributions cannot,
on this data, be separated. The §7 L6 auto-annotation circularity caveat
applies in full to the in-house 56.6\% baseline; the PubMedQA gap
therefore likely overstates the true in-house inflation. The external
Recall@5 is reported here as a directional sanity-check, not as a
head-to-head comparison.

The PubMedQA probe is currently cardiology-only --- the Jaccard
threshold was calibrated for cardiology passages and the filter keyword
list is cardiology-specific. Repeating the probe for gastroenterology
and infectious-disease PubMedQA slices is a follow-up (low cost:
\textasciitilde85 questions × no LLM call, just retrieval + Jaccard);
each new specialty's threshold may need its own calibration because
corpus register varies across specialties.

\subsection{5.10 Regression Check: Cardio / Endo Numbers
Stable}\label{regression-check-cardio-endo-numbers-stable}

The 4-specialist evaluation was specifically designed to preserve the
2-specialist cardio + endo numbers exactly (per-specialty agents have
independent FAISS indices and gold-source annotations; the 4-specialty
eval-script generalisation only changed which keys are iterated, not the
per-key arithmetic). The following cells in §5.3, §5.3.2, §5.4, §5.7,
§5.8 are byte-identical to the prior 2-specialist baseline:

\begin{longtable}[]{@{}
  >{\raggedright\arraybackslash}p{(\linewidth - 6\tabcolsep) * \real{0.2500}}
  >{\raggedright\arraybackslash}p{(\linewidth - 6\tabcolsep) * \real{0.2500}}
  >{\raggedright\arraybackslash}p{(\linewidth - 6\tabcolsep) * \real{0.2500}}
  >{\raggedright\arraybackslash}p{(\linewidth - 6\tabcolsep) * \real{0.2500}}@{}}
\toprule\noalign{}
\begin{minipage}[b]{\linewidth}\raggedright
Metric
\end{minipage} & \begin{minipage}[b]{\linewidth}\raggedright
Cardio
\end{minipage} & \begin{minipage}[b]{\linewidth}\raggedright
Endo
\end{minipage} & \begin{minipage}[b]{\linewidth}\raggedright
Source for prior value
\end{minipage} \\
\midrule\noalign{}
\endhead
\bottomrule\noalign{}
\endlastfoot
Routing accuracy (T1 + T2 + T3) & 100\% (35/35) ✓ & 97.1\% (34/35) ✗
(the 1 miss is endo\_48 routed to infectionist) & §5.1 (original
baseline) \\
Recall@5 T1 (pooled gold) & 59.0\% (23/39) ✓ & 60.6\% (20/33) ✓ & §5.3.2
(BM25 ablation) \\
Recall@5 T2 (pooled gold) & 54.1\% (20/37) ✓ & 52.3\% (23/44) ✓ & §5.3.2
(BM25 ablation) \\
BM25 Recall@5 T1 & 25.6\% (10/39) ✓ & 15.2\% (5/33) ✓ & §5.3.2 (BM25
ablation) \\
BM25 Recall@5 T2 & 43.2\% (16/37) ✓ & 34.1\% (15/44) ✓ & §5.3.2 (BM25
ablation) \\
Multi-judge disagreement set & \{cardio\_40\} ✓ & \{\} ✓ & §5.4 (prior
2-specialty configuration) \\
Minimum-judge faithfulness (cardio T2) & 13/14 = 92.9\% ✓ & --- & §5.4
(prior 2-specialty configuration) \\
§5.3.1 corpus-gap cases & cardio\_35 ✓ (cardio\_23, cardio\_25 retired
from the §5.3.1 list on 2026-05-22 --- both have non-empty
\texttt{gold\_sources} from a later annotation pass and no longer meet
the §5.3.1 empty-gold criterion; underlying corpus-gap claim unchanged)
& endo\_46 ✓ & §5.3.1 (original baseline) \\
\end{longtable}

\textbf{Routing has 1 cell that drifted by exactly 1 case} (endo\_48 →
infectionist instead of endocrinologist). This is \textbf{not a
regression}: the 2-specialty system did not have infectionist as a
routing option; the 4-specialty system does, and \texttt{endo\_48}
happens to be a SGLT2-inhibitor-induced UTI/DKA case where infectionist
is a clinically plausible alternative routing. Every other cardio / endo
cell across the entire §5 evaluation matches the 2-specialty baseline
exactly.

\begin{center}\rule{0.5\linewidth}{0.5pt}\end{center}

\section{6. Discussion}\label{discussion}

Three patterns in the results point to architectural constraints not
visible in the headline numbers.

\subsection{6.1 Precision@K vs.~Hit Rate: The Context-Window Noise
Problem}\label{precisionk-vs.-hit-rate-the-context-window-noise-problem}

MRR@K is the mean reciprocal rank of the first gold-source document
within the top-K retrieved chunks; Recall@K is the fraction of gold
sources retrieved in the top-K window
(\citeproc{ref-manning2008ir}{Manning et al., 2008}). Overall Recall@5
(56.2\% on test) is well below the legacy KeywordHitRate (85.7\% on
test). Recall@5 and Hit Rate are different metrics --- Hit Rate counts a
query as a success if any of the top-K retrieved chunks contains any
expected keyword; Recall@5 counts the fraction of \emph{specific gold
documents} that surface in the top-5 window. The 29-point gap is the
rate at which a single relevant chunk is surfaced alongside four loosely
related ones. With K=5 chunks fed to the generator and only one of them
the gold document, the LLM has to ignore the other four; wider retrieval
windows dilute the gold-document signal: a larger K brings in more
peripheral chunks per query, which increases the chance the LLM
synthesises across loosely related context and fabricates clinical
details not grounded in any single retrieved source. The grid-search
choice K=5, L2 ≤ 1.2 reflects this trade-off (§4.4); re-running the grid
on the new 4-specialty 60-case dev split is a follow-up flagged in §7
L9.

\subsection{6.2 The Nature of Tier 3 Failures: Distance
vs.~Relevance}\label{the-nature-of-tier-3-failures-distance-vs.-relevance}

The original prompt-only ``Insufficient evidence'' fallback failed
completely on Tier 3 (0/29 triggering on the 4-specialty test split).
FAISS returns K results regardless of absolute relevance; the LLM treats
whatever is retrieved as context and writes from it; the in-prompt
fallback instruction loses to the LLM's training to be helpful. The L2
distance threshold filters by chunk quality but does not act as a
semantic relevance gate. A numeric pre-LLM refusal gate was added (§5.5)
that checks \texttt{min(L2)\ \textgreater{}\ L2\_REJECT\_MIN} before the
generation call. On the 2-specialty data this raised Tier 3 refusal from
0/15 to 12/15 (80\%) at the cost of a 49\% T1/T2 FP rate; on the new
4-specialty data the trade-off curve has shifted unfavourably (gastro /
infect μ\_corpus ≈ 0.88 vs cardio μ ≈ 0.87 and endo μ ≈ 0.89 --- but the
new corpora's T1/T2 min-L2 range extends to 1.107, vs 2-spec's 1.07),
and no threshold simultaneously meets both targets (§5.5). The re-tuned
\texttt{L2\_REJECT\_MIN\ =\ 1.020} yields 20.7\% T3 recall and 17.1\%
T1/T2 FP --- a more permissive gate that lets through most legitimate
queries but catches only a fifth of true T3 cases. \textbf{Reliable
out-of-scope detection on this corpus requires a signal in addition to
top-K L2 distance}: a two-stage gate combining the numeric pre-filter
with an LLM-as-classifier confirmer is the architecture now needed.

\subsection{6.3 Epistemic Bounds of Same-Family
Evaluation}\label{epistemic-bounds-of-same-family-evaluation}

The two Yandex judges disagree on exactly one test-split case ---
\texttt{cardio\_40} (Tier 2 cardiology) --- and that single disagreement
makes it possible to characterise the kind of call the primary YandexGPT
judge accepts and the secondary YandexGPT-Lite judge rejects. The query
asks for the likely diagnosis of a 30-year-old male with resuscitated
out-of-hospital cardiac arrest, prolonged QTc of 510 ms, and a sister
who had a similar event at age 25 --- a presentation that strongly
suggests congenital long QT syndrome. The retrieved context contains a
tangentially related case (30-something woman with new-onset seizure
activity and prolonged QTc 500--530 ms leading to Torsades de Pointes)
which explicitly attributes the prolongation to herbal-remedy-induced
\emph{acquired} LQTS while noting that ``normal QTc does not exclude
congenital LQTS.'' The generated answer paraphrases this related case,
then infers congenital LQTS for the new patient citing the family
history. The primary judge accepts this as a faithful paraphrase plus
logical inference allowed by the rules and returns \texttt{FAITHFUL}.
The secondary judge rejects it as introducing a specific diagnosis
(congenital LQTS) not directly named in the retrieved context and
returns \texttt{HALLUCINATION}.

The flagship YandexGPT primary judge accepts inferences from
related-but-distinct context; the smaller YandexGPT-Lite secondary judge
requires the specific diagnosis label to appear in the retrieved tokens
before returning FAITHFUL. The secondary judge applies a stricter
standard --- direct naming in retrieved tokens rather than inference ---
and the minimum-rate protocol adopts that standard deliberately,
yielding a conservative faithfulness lower bound: under the stricter
Yandex judge the test-split 2-Yandex-judge minimum-judge rate is 131/132
= 99.2\% {[}Wilson 95\% CI 95.8\%--99.9\%{]}, with the single
Yandex-pair disagreement remaining on \texttt{cardio\_40} --- exactly as
in the prior 2-specialty configuration. The denominator grew from 70 to
132 (after excluding 7 Tier-3 fallback cases and 1 None-verdict case
from the 140-case test split). This 99.2\% rate is the \emph{upper}
bound on faithfulness, valid under same-family bias.

\textbf{The lower bound --- and the bound to quote against external
evaluations --- is the 3-judge cross-vendor rate of 127/132 = 96.2\%
{[}91.4\%--98.4\%{]}.} The cross-vendor tertiary judge (OpenAI
GPT-OSS-120B, configured 2026-05-23) flagged four additional cases as
HALLUCINATION that both Yandex judges marked FAITHFUL:
\texttt{cardio\_17}, \texttt{endo\_21}, \texttt{endo\_38},
\texttt{infect\_34}. These four are the empirical fingerprint of the
same-family blind spot Zheng et al.
(\citeproc{ref-zheng2023mtbench}{2023}) characterise: a model from a
different vendor family caught hallucinations that two judges from the
same vendor as the generator did not. The Yandex-pair
\texttt{cardio\_40} disagreement remains historically interesting as a
clinical-reasoning case study, but the headline-shaping disagreements
after adding the cross-vendor judge are the four cross-vendor flags.
Each warrants individual chart review; the §5.4 disagreement table
catalogues them.

The same-family secondary judge (YandexGPT-Lite alongside YandexGPT) was
a budget-driven choice rather than a methodological preference: the API
access available within project scope was Yandex's, and hiring a
clinical specialist for manual chart review was out of reach. The
cross-vendor third judge in §5.4 was added specifically to bound the
same-family-bias risk that the two-Yandex axis alone cannot rule out.

\begin{center}\rule{0.5\linewidth}{0.5pt}\end{center}

\section{7. Limitations}\label{limitations}

\begin{enumerate}
\def\labelenumi{\arabic{enumi}.}
\item
  \textbf{Golden dataset size.} The evaluation uses 200 cases (50 per
  specialty × 4 specialties). The doubling from 100 → 200 tightens every
  Wilson lower bound by 5--10 pp; an even larger test set (1,000+ cases)
  would tighten further and expose rarer failure modes, especially on
  tier-3 refusal where the per-tier n = 29 still yields wide CIs.
\item
  \textbf{Cardiology + new-specialty corpus coverage gaps --- concrete
  case list.} §5.3.1 currently documents six Tier 1/2 cases where the
  gold-source annotator could not find any document with ≥1
  expected-keyword match in the top-20 retrieval window:
  \texttt{cardio\_35} and \texttt{endo\_46} (original baseline);
  \texttt{gastro\_39}, \texttt{gastro\_44}, \texttt{infect\_21},
  \texttt{infect\_39} (added with the four-specialty expansion). On
  2026-05-22, \texttt{gastro\_37} and \texttt{infect\_14} were retired
  from the audit (keyword-set mismatches, not corpus gaps --- see §5.3.1
  retirement note); \texttt{cardio\_23} and \texttt{cardio\_25} were
  retired in the same pass (non-empty \texttt{gold\_sources} from a
  later annotation pass). The remaining six are characterised in §5.3.1;
  the underlying-source-material remediation list is concrete
  (procedural cardiology guidelines for \texttt{cardio\_35}; PJI
  references for \texttt{infect\_39}; metabolic-liver-disease references
  for \texttt{gastro\_39}), and is the next ingestion step.
\item
  \textbf{Token limit constraints.} The Yandex embedding model has a
  hard limit of 2,048 tokens per request. Approximately 20 chunks in the
  endocrinology corpus required automatic truncation during index
  building. A load-time mean-word-length filter was added to the index
  builder (\texttt{\_MAX\_MEAN\_WORD\_LEN\_CHARS\ =\ 15}) that drops
  PDF-extraction-artifact chunks with concatenated text (no inter-word
  spaces, mean word length 30--170 chars instead of normal English's
  \textasciitilde6--7) before they reach the embedding model. The filter
  drops 354 / 9,024 = 3.9\% of gastro chunks and 236 / 7,712 = 3.1\% of
  infect chunks (almost entirely author-affiliation blocks of
  multi-author papers); 0 / 7,730 cardio chunks and 0 / 37,791 endo
  chunks are affected.
\item
  \textbf{Single-language corpus.} All source documents are in English.
  The system has not been validated for multilingual queries or
  non-English medical literature; the adversarial test (§5.2.1) probes
  Russian / French / Spanish surface vocabulary mapped against the
  English corpus, with 16/16 = 100\% routing accuracy, but the retrieval
  and generation are not validated on non-English queries.
\item
  \textbf{LLM-as-a-judge circularity (cross-vendor judge added
  2026-05-23).} Faithfulness is evaluated by three judges given the
  identical strict prompt: two from different Yandex model families
  (\texttt{yandexgpt/latest} and \texttt{yandexgpt-lite/latest}) and one
  cross-vendor (\texttt{openai/gpt-oss-120b:free} via OpenRouter). κ and
  Gwet's AC1 values for all three judge pairs are reported in §5.4. The
  minimum-judge rate --- a case is counted FAITHFUL only if every
  configured judge agrees --- is the headline number used in §8.
  Residual risk after this stage: only one vendor family outside Yandex
  is sampled; a second cross-vendor judge (e.g.~an Anthropic or Google
  model) would further shrink the residual single-vendor-blind-spot
  risk, but the marginal benefit beyond the first cross-vendor judge is
  bounded by Zheng et al.'s (\citeproc{ref-zheng2023mtbench}{Zheng et
  al., 2023}) same-family inflation estimate (5--25 pp), of which this
  run measured 3 pp on this specific \{Yandex, OpenAI\} pair.
\item
  \textbf{Auto-annotation circularity in Recall@K.} The 162 of 168 Tier
  1/2 gold\_sources annotations used as the Recall@K denominator were
  produced by the gold-source annotator in \texttt{-\/-auto} mode, which
  selected up to three documents per case from the top-20 retrieval
  output of the same FAISS+embedding system being evaluated. Recall@K
  therefore measures `fraction of keyword-positive top-20 documents that
  surface in the top-5 window', not `fraction of ground-truth answer
  documents retrieved'. An independent annotation pass (human-curated
  against the full corpus, or against a benchmark like PubMedQA --- see
  §5.9) would break the circularity; this is the natural next step for
  any future work that intends to compare against external retrieval
  baselines.
\item
  \textbf{Refusal gate.} The single-threshold numeric gate cannot
  simultaneously meet both targets (≥80\% T3 recall, ≤5\% T1/T2 FP) at
  4-specialty scale (see §5.5 for the threshold sweep, overlap-zone
  characterisation, and chosen fallback). The proposed two-stage gate
  (per-specialty numeric pre-filter + LLM confirmer on the overlap zone)
  is the natural next experiment; this evaluation does not implement it.
  Writing this limitation required accepting the L2 threshold sweep
  result rather than re-running with a wider rejection band: every
  widening that improved Tier 3 catch came at the cost of false
  rejections on Tier 1/2 cases, leaving the original target unmet rather
  than apparently met.
\item
  \textbf{Earlier multijudge reconciliation is on n = 70.} A prior
  reconciliation pass characterises the non-determinism of
  \texttt{yandexgpt-lite/latest} at temperature=0 on the 2-specialty
  test split (n = 70). A 4-specialty (n = 140) reconciliation re-run
  would compare the earlier disagreement set (just cardio\_40) against a
  fresh re-run to characterise variance under the larger sample --- a
  follow-up.
\item
  \textbf{Retrieval hyperparameter grid pre-dates the four-specialty
  expansion.} The \texttt{K\ ×\ L2} grid (§4.4) and the chunk-size grid
  (§4.3) were both run on the original 30-case cardio + endo dev split.
  Whether \texttt{K=5,\ L2\ ≤\ 1.2} and
  \texttt{CHUNK\_SIZE\_WORDS\ =\ 400} remain Pareto-optimal on the
  60-case 4-specialty dev split is an open question. The cost is one
  re-run of the K × L2 grid sweep (no API calls beyond the existing
  FAISS retrieval) plus one optional re-build of the FAISS indices if a
  different chunk size were selected (≈\$30--40 + \textasciitilde5 h on
  the Yandex Embedding API). The TF-IDF router question that originally
  lived here --- the TF-IDF model was 2-class until 2026-05-22 --- was
  \textbf{resolved} by retraining on the 60-case 4-spec dev split
  (one-line script change, \texttt{random\_state=42} pinned); see §5.1
  for the post-retrain numbers.
\end{enumerate}

\begin{center}\rule{0.5\linewidth}{0.5pt}\end{center}

\section{8. Conclusion}\label{conclusion}

\subsection{Research question closure}\label{research-question-closure}

The evaluation addresses the three architectural questions posed in §1:

\emph{RQ1 --- Does an LLM-based query router add measurable clinical
value over a deterministic keyword-matching baseline?} Yes: the LLM
router achieves 95.7\% (134/140) routing accuracy vs 79.3\% (111/140)
for the keyword baseline on the n = 140 test split --- a 16.4 pp lead
that holds across all four specialties and is most pronounced on
cross-domain cases.

\emph{RQ2 --- How does retrieval quality degrade across difficulty
tiers?} Retrieval quality holds roughly constant across Tier 1 (pooled
78/140 = 55.7\% Recall@5) and Tier 2 (pooled 84/148 = 56.8\%) on the
held-out test split. The design hypothesis that Tier 2 (peripheral)
queries would underperform Tier 1 (core) queries is not supported by
these data; per-specialty variation dominates per-tier variation, with
the strongest per-cell rate (62.5\%, T2 Infectiology) and the weakest
(50.0\%, T1 Infectiology) separated by 12.5 pp.

\emph{RQ3 --- Can a multi-judge LLM-as-judge protocol reliably detect
hallucinations in generated medical advice?} Within bounds: the
minimum-judge rule across two same-family Yandex judges yields a 99.2\%
upper-bound faithfulness rate, with a 3-judge cross-vendor lower bound
of 96.2\% on the test split. The measured −3.0 pp cross-vendor signal
from the third OpenAI judge is consistent with the low end of the 5--25
pp same-family inflation range Zheng et al.
(\citeproc{ref-zheng2023mtbench}{2023}) characterise.

\subsection{Principal architectural
finding}\label{principal-architectural-finding}

The clearest architectural conclusion from this evaluation is that
single-threshold L2-distance gating is insufficient for clinical-grade
out-of-scope refusal at four-specialty scale. The system meets or
approaches its other evaluation targets --- 95.7\% routing accuracy,
96.2\%/99.2\% bounded faithfulness, 56.2\% Recall@5 --- but the 20.7\%
Tier 3 catch on the n = 140 held-out test split sits below any threshold
that would justify clinical use, and the threshold sweep (§5.5) shows no
scalar \texttt{L2\_REJECT\_MIN} value simultaneously satisfies both
target rates.

Headline metrics are reported on the 140-case held-out test split
(§5.8), which excludes the 60 development cases (15 per specialty × 4
specialties) used to tune K, L2 threshold, chunk size, and the
refusal-gate threshold. Faithfulness is reported under the minimum-judge
rule (a case counts FAITHFUL only if every configured judge agrees). The
numeric refusal gate (§5.5) replaces the prompt-only fallback that
previously failed on every Tier 3 case; the gate threshold was re-tuned
for the 4-specialty corpus. \textbf{All headline numbers in this
Conclusion are computed on the held-out test split (n = 140); the
full-set numbers (n = 200) are presented separately in §5.7.} The
headline metrics on the held-out test split:

\begin{itemize}
\tightlist
\item
  \textbf{Routing accuracy on the held-out test split: 95.7\% (134/140)
  {[}Wilson 95\% CI 91.0\%--98.0\%{]}} (§5.1, §5.8). Cardiology 100\%
  (35/35), Endocrinology 97.1\% (34/35), Gastroenterology 88.6\%
  (31/35), Infectiology 97.1\% (34/35). The 6 LLM misses are all
  defensible cross-specialty ambiguities (e.g.~\texttt{gastro\_37} viral
  gastroenteritis → infectionist; \texttt{endo\_48} SGLT2-inhibitor UTI
  complication → infectionist), not bugs. On the 14 ambiguous
  cross-domain cases (§5.2) the LLM stays inside \texttt{valid\_domains}
  on 14/14 (after the 2026-05-22 widening of
  \texttt{ambig\_8.valid\_domains} to admit the clinically defensible
  gastroenterology routing); on the 64-case adversarial set (§5.2.1) it
  scores 98.4\% (63/64) overall on the historical eval and would lift to
  64/64 = 100\% on the next pass after the parallel
  \texttt{adv\_amb\_6.valid\_domains} widening, with 100\% on the
  misspelled, non-English, and dominant-pathology-mismatch categories.
\item
  \textbf{Retrieval Recall@5 on the held-out test split: 56.2\%
  (162/288) {[}Wilson 95\% CI 50.5\%--61.9\%{]}} (§5.3, §5.8) across 4
  specialties. Per-specialty: Cardiology 56.6\% (43/76), Endocrinology
  55.8\% (43/77), Gastroenterology 56.3\% (40/71), Infectiology 56.3\%
  (36/64). Cardio and endo cells are byte-identical to the prior
  2-specialty baseline --- a verified regression check. FAISS dense
  retrieval outperforms BM25 by 22.2 pp overall (56.2\% vs 34.0\%); BM25
  narrows the gap on Tier 2 cases where exact-entity-name matches help.
  The legacy KeywordHitRate (85.7\% on test) registers any keyword
  co-occurrence in the top-5 window and remains a loose secondary
  signal.
\item
  \textbf{Out-of-scope refusal on the held-out test split: 6/29 = 20.7\%
  {[}Wilson 95\% CI 9.8\%--38.4\%{]}} (§5.5) at the re-tuned
  \texttt{L2\_REJECT\_MIN\ =\ 1.020}. T1/T2 FP rate: 19/111 = 17.1\%
  {[}11.2\%--25.2\%{]}. The single-threshold gate cannot simultaneously
  satisfy the ≥80\% T3 recall and ≤5\% T1/T2 FP targets at 4-specialty
  scale because the in-corpus / out-of-corpus min-L2 distributions
  overlap heavily (gastro and infect μ\_corpus ≈ 0.88, with the T3 mean
  only 0.10 above) --- a two-stage gate (numeric pre-filter +
  LLM-as-classifier confirmer) is now urgent (§7 L7).
\item
  \textbf{Minimum-judge faithfulness on the held-out test split: 127/132
  = 96.2\% {[}Wilson 95\% CI 91.4\%--98.4\%{]}} with a 3-judge
  cross-vendor panel (§5.4). The same-family 2-Yandex-judge upper bound
  is 131/132 = 99.2\% {[}95.8\%--99.9\%{]}; the cross-vendor tertiary
  judge (OpenAI GPT-OSS-120B, configured 2026-05-23) added four
  HALLUCINATION verdicts on cases both Yandex judges marked FAITHFUL ---
  \texttt{cardio\_17}, \texttt{endo\_21}, \texttt{endo\_38},
  \texttt{infect\_34} --- for a −3.0 pp drop that is the empirical
  fingerprint of the same-family bias Zheng et al.
  (\citeproc{ref-zheng2023mtbench}{2023}) predict. The Yandex pair's
  single historical disagreement (\texttt{cardio\_40}, Tier 2 cardiology
  --- congenital LQTS) is unchanged from the prior 2-specialty
  configuration. The 91.4\% Wilson lower bound on the 3-judge rate is
  the number to quote against external LLM-as-judge faithfulness
  results; the 2-Yandex-judge rate is retained as the upper bound under
  same-family bias.
\end{itemize}

The hyperparameter grid search (K × L2 threshold) ran on the 30-case
original development split (§4.4) and selected K=5, L2 ≤ 1.2; re-running
it on the 60-case 4-specialty dev split is a follow-up. The chunk size
choice (400 words) and keyword-stripping were validated only under joint
application (§4.5). All four specialist agents --- cardiologist,
endocrinologist, gastroenterologist, infectionist --- are evaluated
end-to-end. \textbf{Adding any further specialty is a single registry
entry plus a corpus and FAISS+BM25 index build}; no code changes to the
evaluation pipeline are required (every per-domain loop was generalised
to iterate over the registered specialties rather than hardcoded
2-specialty tuples).

Closing the 20.7\% gap toward clinical-deployment refusal rates remains
the open problem that gates this prototype's transition from academic
evaluation to clinical use; the two-stage L2-plus-classifier remediation
described in §7 L7 is the most direct architectural path.

\begin{center}\rule{0.5\linewidth}{0.5pt}\end{center}

\section{Appendix A}\label{appendix-a}

Detailed evaluation tables routed from the main body for readability.

\subsection{A.1 --- Cross-domain ambiguous routing, 14
cases}\label{a.1-cross-domain-ambiguous-routing-14-cases}

Full per-case routing decisions from §5.2. The LLM router's prediction
is compared against the keyword and 4-class TF-IDF baselines.

\textbf{Table A.1.} Ambiguous-case routing decisions across all three
routers.

\begin{longtable}[]{@{}
  >{\raggedright\arraybackslash}p{(\linewidth - 10\tabcolsep) * \real{0.1667}}
  >{\raggedright\arraybackslash}p{(\linewidth - 10\tabcolsep) * \real{0.1667}}
  >{\raggedright\arraybackslash}p{(\linewidth - 10\tabcolsep) * \real{0.1667}}
  >{\raggedright\arraybackslash}p{(\linewidth - 10\tabcolsep) * \real{0.1667}}
  >{\raggedright\arraybackslash}p{(\linewidth - 10\tabcolsep) * \real{0.1667}}
  >{\raggedright\arraybackslash}p{(\linewidth - 10\tabcolsep) * \real{0.1667}}@{}}
\toprule\noalign{}
\begin{minipage}[b]{\linewidth}\raggedright
ID
\end{minipage} & \begin{minipage}[b]{\linewidth}\raggedright
Clinical Scenario
\end{minipage} & \begin{minipage}[b]{\linewidth}\raggedright
LLM
\end{minipage} & \begin{minipage}[b]{\linewidth}\raggedright
Keyword
\end{minipage} & \begin{minipage}[b]{\linewidth}\raggedright
TF-IDF
\end{minipage} & \begin{minipage}[b]{\linewidth}\raggedright
Valid Domains
\end{minipage} \\
\midrule\noalign{}
\endhead
\bottomrule\noalign{}
\endlastfoot
ambig\_1 & Diabetic cardiomyopathy (HbA1c 9.2\%, EF 40\%) & cardiologist
& cardiologist & cardiologist & cardio, endo \\
ambig\_2 & Thyroid-induced atrial fibrillation (Graves', HR 130) &
endocrinologist & cardiologist & endocrinologist & cardio, endo \\
ambig\_3 & SGLT2 inhibitor cardioprotection in ACS & cardiologist &
cardiologist & cardiologist & cardio, endo \\
ambig\_4 & Hyperaldosteronism with resistant hypertension (K+ 2.9) &
endocrinologist & cardiologist & endocrinologist & cardio, endo \\
ambig\_5 & Pheochromocytoma with Takotsubo cardiomyopathy (BP 240/140) &
endocrinologist & cardiologist & cardiologist & cardio, endo \\
ambig\_6 & Amiodarone-induced hypothyroidism (TSH 45) & endocrinologist
& cardiologist & endocrinologist & cardio, endo \\
ambig\_7 & Metabolic syndrome with exertional angina (BMI 38, +stress
test) & cardiologist & cardiologist & endocrinologist & cardio, endo \\
ambig\_8 & Carcinoid heart disease (right-sided valve lesions, ↑5-HIAA)
& gastroenterologist & cardiologist & endocrinologist & cardio, endo,
gastro \\
ambig\_9 & H. pylori peptic ulcer with iron-deficiency anaemia &
gastroenterologist & gastroenterologist & gastroenterologist & gastro,
endo \\
ambig\_10 & Sepsis with new-onset hyperglycaemia & endocrinologist &
endocrinologist & gastroenterologist & infect, endo \\
ambig\_11 & HIV-positive patient with dilated cardiomyopathy &
cardiologist & infectionist & cardiologist & infect, cardio \\
ambig\_12 & Autoimmune hepatitis with thyroid disease &
gastroenterologist & endocrinologist & endocrinologist & gastro, endo \\
ambig\_13 & C. difficile colitis post-antibiotic & infectionist &
gastroenterologist & infectionist & infect, gastro \\
ambig\_14 & Liver cirrhosis with spontaneous bacterial peritonitis &
infectionist & gastroenterologist & gastroenterologist & gastro,
infect \\
\end{longtable}

\subsection{A.2 --- Retriever comparison, full per-Tier per-specialty
breakdown}\label{a.2-retriever-comparison-full-per-tier-per-specialty-breakdown}

Per-Tier per-specialty Recall@5 and MRR@5 for FAISS dense, BM25 sparse,
random baseline, and oracle ceiling on the held-out 140-case test split.
Pooled gold-doc Bernoulli with Wilson 95\% CI on Recall@5; MRR@5 95\%
CIs are percentile-method bootstrap intervals (B=10000, RNG seed=12345).

\textbf{Table A.2.} FAISS vs BM25 vs Random vs Oracle, stratified by
specialty and Tier.

\begin{longtable}[]{@{}
  >{\raggedright\arraybackslash}p{(\linewidth - 14\tabcolsep) * \real{0.1250}}
  >{\raggedright\arraybackslash}p{(\linewidth - 14\tabcolsep) * \real{0.1250}}
  >{\raggedright\arraybackslash}p{(\linewidth - 14\tabcolsep) * \real{0.1250}}
  >{\raggedright\arraybackslash}p{(\linewidth - 14\tabcolsep) * \real{0.1250}}
  >{\raggedright\arraybackslash}p{(\linewidth - 14\tabcolsep) * \real{0.1250}}
  >{\raggedright\arraybackslash}p{(\linewidth - 14\tabcolsep) * \real{0.1250}}
  >{\raggedright\arraybackslash}p{(\linewidth - 14\tabcolsep) * \real{0.1250}}
  >{\raggedright\arraybackslash}p{(\linewidth - 14\tabcolsep) * \real{0.1250}}@{}}
\toprule\noalign{}
\begin{minipage}[b]{\linewidth}\raggedright
Domain
\end{minipage} & \begin{minipage}[b]{\linewidth}\raggedright
Tier
\end{minipage} & \begin{minipage}[b]{\linewidth}\raggedright
FAISS Recall@5
\end{minipage} & \begin{minipage}[b]{\linewidth}\raggedright
BM25 Recall@5
\end{minipage} & \begin{minipage}[b]{\linewidth}\raggedright
Random Recall@5
\end{minipage} & \begin{minipage}[b]{\linewidth}\raggedright
Oracle Recall@5
\end{minipage} & \begin{minipage}[b]{\linewidth}\raggedright
FAISS MRR@5 {[}95\% CI{]}
\end{minipage} & \begin{minipage}[b]{\linewidth}\raggedright
BM25 MRR@5 {[}95\% CI{]}
\end{minipage} \\
\midrule\noalign{}
\endhead
\bottomrule\noalign{}
\endlastfoot
Cardiology & T1 (core) & 59.0\% (23/39) {[}43.4\%--72.9\%{]} & 25.6\%
(10/39) {[}14.6\%--41.1\%{]} & 2.6\% (1/39) {[}0.5\%--13.2\%{]} & 100\%
(39/39) {[}91.0\%--100\%{]} & 0.737 {[}0.577--0.885{]} & 0.442
{[}0.212--0.673{]} \\
Cardiology & T2 & 54.1\% (20/37) {[}38.4\%--69.0\%{]} & 43.2\% (16/37)
{[}28.7\%--59.1\%{]} & 13.5\% (5/37) {[}5.9\%--28.0\%{]} & 100\% (37/37)
{[}90.6\%--100\%{]} & 0.567 {[}0.345--0.785{]} & 0.562
{[}0.338--0.785{]} \\
Endocrinology & T1 & 60.6\% (20/33) {[}43.7\%--75.3\%{]} & 15.2\% (5/33)
{[}6.7\%--30.9\%{]} & 0.0\% (0/33) {[}0.0\%--10.4\%{]} & 100\% (33/33)
{[}89.6\%--100\%{]} & 0.773 {[}0.591--0.939{]} & 0.348
{[}0.106--0.621{]} \\
Endocrinology & T2 & 52.3\% (23/44) {[}37.9\%--66.2\%{]} & 34.1\%
(15/44) {[}21.9\%--48.9\%{]} & 0.0\% (0/44) {[}0.0\%--8.0\%{]} & 100\%
(44/44) {[}92.0\%--100\%{]} & 0.677 {[}0.458--0.875{]} & 0.492
{[}0.285--0.700{]} \\
Gastroenterology & T1 & 52.8\% (19/36) {[}37.0\%--68.0\%{]} & 41.7\%
(15/36) {[}27.1\%--57.8\%{]} & 5.6\% (2/36) {[}1.5\%--18.1\%{]} & 100\%
(36/36) {[}90.4\%--100\%{]} & 0.654 {[}0.436--0.859{]} & 0.513
{[}0.282--0.744{]} \\
Gastroenterology & T2 & 60.0\% (21/35) {[}43.6\%--74.4\%{]} & 37.1\%
(13/35) {[}23.2\%--53.7\%{]} & 2.9\% (1/35) {[}0.5\%--14.5\%{]} & 100\%
(35/35) {[}90.1\%--100\%{]} & 0.808 {[}0.615--0.962{]} & 0.596
{[}0.346--0.846{]} \\
Infectiology & T1 & 50.0\% (16/32) {[}33.6\%--66.4\%{]} & 31.2\% (10/32)
{[}18.0\%--48.6\%{]} & 0.0\% (0/32) {[}0.0\%--10.7\%{]} & 100\% (32/32)
{[}89.3\%--100\%{]} & 0.750 {[}0.528--0.944{]} & 0.454
{[}0.188--0.708{]} \\
Infectiology & T2 & 62.5\% (20/32) {[}45.3\%--77.1\%{]} & 43.8\% (14/32)
{[}28.2\%--60.7\%{]} & 0.0\% (0/32) {[}0.0\%--10.7\%{]} & 100\% (32/32)
{[}89.3\%--100\%{]} & 0.642 {[}0.419--0.857{]} & 0.667
{[}0.452--0.881{]} \\
\textbf{Overall (T1+T2)} & --- & \textbf{56.2\% (162/288)
{[}50.5\%--61.9\%{]}} & \textbf{34.0\% (98/288) {[}28.8\%--39.7\%{]}} &
\textbf{3.8\% (11/288) {[}2.1\%--6.7\%{]}} & \textbf{100\% (288/288)
{[}98.7\%--100\%{]}} & \textbf{0.698 {[}0.624--0.770{]}} & \textbf{0.514
{[}0.429--0.597{]}} \\
\end{longtable}

\subsection{A.3 --- Judge-pair agreement (Cohen's κ and Gwet's
AC1)}\label{a.3-judge-pair-agreement-cohens-ux3ba-and-gwets-ac1}

Pairwise inter-rater agreement across the three judges on the n = 132
intersection (test split excluding fallback-skipped and None-verdict
cases). Cohen's κ degenerates when one rater's marginal class
probability is 0; Gwet's AC1 (\citeproc{ref-gwet2008ac1}{Gwet, 2008})
remains interpretable in that regime.

\textbf{Table A.3.} Judge-pair agreement.

\begin{longtable}[]{@{}
  >{\raggedright\arraybackslash}p{(\linewidth - 10\tabcolsep) * \real{0.1364}}
  >{\raggedleft\arraybackslash}p{(\linewidth - 10\tabcolsep) * \real{0.1818}}
  >{\raggedleft\arraybackslash}p{(\linewidth - 10\tabcolsep) * \real{0.1818}}
  >{\raggedleft\arraybackslash}p{(\linewidth - 10\tabcolsep) * \real{0.1818}}
  >{\raggedleft\arraybackslash}p{(\linewidth - 10\tabcolsep) * \real{0.1818}}
  >{\raggedright\arraybackslash}p{(\linewidth - 10\tabcolsep) * \real{0.1364}}@{}}
\toprule\noalign{}
\begin{minipage}[b]{\linewidth}\raggedright
Pair
\end{minipage} & \begin{minipage}[b]{\linewidth}\raggedleft
n
\end{minipage} & \begin{minipage}[b]{\linewidth}\raggedleft
Agreements
\end{minipage} & \begin{minipage}[b]{\linewidth}\raggedleft
Cohen's κ
\end{minipage} & \begin{minipage}[b]{\linewidth}\raggedleft
Gwet's AC1
\end{minipage} & \begin{minipage}[b]{\linewidth}\raggedright
Landis \& Koch
\end{minipage} \\
\midrule\noalign{}
\endhead
\bottomrule\noalign{}
\endlastfoot
(primary, secondary) & 132 & 131 & 0.000 & 0.992 & degenerate κ; via
AC1: almost perfect \\
(primary, tertiary) & 132 & 128 & 0.000 & 0.969 & degenerate κ; via AC1:
almost perfect \\
(secondary, tertiary) & 132 & 127 & −0.012 & 0.961 & less than chance
(κ); almost perfect (AC1) \\
\end{longtable}

\subsection{A.4 --- Refusal-gate precision/recall, full per-Tier
per-specialty
breakdown}\label{a.4-refusal-gate-precisionrecall-full-per-tier-per-specialty-breakdown}

Per-cell refusal rates at the chosen threshold
\texttt{L2\_REJECT\_MIN\ =\ 1.020} on the 140-case held-out test split.
The overall Tier 3 (positive-class) and Tier 1/2 (negative-class FP)
rows are reproduced in §5.5 main; the 12 per-Tier per-specialty cells
below are the full breakdown.

\textbf{Table A.4.} Per-Tier per-specialty refusal rate.

\begin{longtable}[]{@{}llll@{}}
\toprule\noalign{}
Stratum & Cases & Refused by gate & Refusal rate \\
\midrule\noalign{}
\endhead
\bottomrule\noalign{}
\endlastfoot
T1 Cardiology & 13 & 1 & 7.7\% \\
T2 Cardiology & 14 & 2 & 14.3\% \\
T3 Cardiology & 8 & 0 & 0.0\% \\
T1 Endocrinology & 12 & 0 & 0.0\% \\
T2 Endocrinology & 16 & 0 & 0.0\% \\
T3 Endocrinology & 7 & 0 & 0.0\% \\
T1 Gastroenterology & 13 & 2 & 15.4\% \\
T2 Gastroenterology & 15 & 1 & 6.7\% \\
T3 Gastroenterology & 7 & 3 & 42.9\% \\
T1 Infectiology & 13 & 3 & 23.1\% \\
T2 Infectiology & 15 & 10 & 66.7\% \\
T3 Infectiology & 7 & 3 & 42.9\% \\
\end{longtable}

References

\protect\phantomsection\label{refs}
\begin{CSLReferences}{1}{0}
\bibitem[\citeproctext]{ref-es2023ragas}
Es, S., James, J., Espinosa-Anke, L., \& Schockaert, S. (2023).
\emph{{RAGAS}: Automated evaluation of retrieval augmented generation}.
\url{https://arxiv.org/abs/2309.15217}

\bibitem[\citeproctext]{ref-gwet2008ac1}
Gwet, K. L. (2008). Computing inter-rater reliability and its variance
in the presence of high agreement. \emph{British Journal of Mathematical
and Statistical Psychology}, \emph{61}(1), 29--48.
\url{https://doi.org/10.1348/000711006X126600}

\bibitem[\citeproctext]{ref-jin2019pubmedqa}
Jin, Q., Dhingra, B., Liu, Z., Cohen, W. W., \& Lu, X. (2019).
{PubMedQA}: A dataset for biomedical research question answering.
\emph{Proceedings of the 2019 Conference on Empirical Methods in Natural
Language Processing and the 9th International Joint Conference on
Natural Language Processing (EMNLP-IJCNLP)}, 2567--2577.
\url{https://arxiv.org/abs/1909.06146}

\bibitem[\citeproctext]{ref-karpukhin2020dpr}
Karpukhin, V., Oğuz, B., Min, S., Lewis, P., Wu, L., Edunov, S., Chen,
D., \& Yih, W. (2020). Dense passage retrieval for open-domain question
answering. \emph{Proceedings of the 2020 Conference on Empirical Methods
in Natural Language Processing (EMNLP)}, 6769--6781.
\url{https://arxiv.org/abs/2004.04906}

\bibitem[\citeproctext]{ref-kim2024mdagents}
Kim, Y., Park, C., Jeong, H., Chan, Y. S., Xu, X., McDuff, D., Breazeal,
C., \& Park, H. W. (2024). {MDAgents}: An adaptive collaboration of
{LLM}s for medical decision making. \emph{Advances in Neural Information
Processing Systems}. \url{https://arxiv.org/abs/2404.15155}

\bibitem[\citeproctext]{ref-lewis2020rag}
Lewis, P., Perez, E., Piktus, A., Petroni, F., Karpukhin, V., Goyal, N.,
Küttler, H., Lewis, M., Yih, W., Rocktäschel, T., Riedel, S., \& Kiela,
D. (2020). Retrieval-augmented generation for knowledge-intensive {NLP}
tasks. \emph{Advances in Neural Information Processing Systems},
\emph{33}. \url{https://arxiv.org/abs/2005.11401}

\bibitem[\citeproctext]{ref-manning2008ir}
Manning, C. D., Raghavan, P., \& Schütze, H. (2008). \emph{Introduction
to information retrieval}. Cambridge University Press.

\bibitem[\citeproctext]{ref-robertson2009bm25}
Robertson, S., \& Zaragoza, H. (2009). The probabilistic relevance
framework: {BM25} and beyond. \emph{Foundations and Trends in
Information Retrieval}, \emph{3}(4), 333--389.
\url{https://doi.org/10.1561/1500000019}

\bibitem[\citeproctext]{ref-singhal2023medpalm}
Singhal, K., Azizi, S., Tu, T., Mahdavi, S. S., Wei, J., Chung, H. W.,
Scales, N., Tanwani, A., Cole-Lewis, H., Pfohl, S., Payne, P.,
Seneviratne, M., Gamble, P., Kelly, C., Babiker, A., Schärli, N.,
Chowdhery, A., Mansfield, P., Demner-Fushman, D., \ldots{} Natarajan, V.
(2023). Large language models encode clinical knowledge. \emph{Nature},
\emph{620}(7972), 172--180.
\url{https://doi.org/10.1038/s41586-023-06291-2}

\bibitem[\citeproctext]{ref-tsatsaronis2015bioasq}
Tsatsaronis, G., Balikas, G., Malakasiotis, P., Partalas, I., Zschunke,
M., Alvers, M. R., Weissenborn, D., Krithara, A., Petridis, S.,
Polychronopoulos, D., Almirantis, Y., Pavlopoulos, J., Baskiotis, N.,
Gallinari, P., Artières, T., Ngonga Ngomo, A.-C., Heino, N., Gaussier,
E., Barrio-Alvers, L., \ldots{} Paliouras, G. (2015). An overview of the
{BioASQ} large-scale biomedical semantic indexing and question answering
competition. \emph{BMC Bioinformatics}, \emph{16}(1), 138.
\url{https://doi.org/10.1186/s12859-015-0564-6}

\bibitem[\citeproctext]{ref-xiong2024medrag}
Xiong, G., Jin, Q., Lu, Z., \& Zhang, A. (2024). \emph{Benchmarking
retrieval-augmented generation for medicine}.
\url{https://arxiv.org/abs/2402.13178}

\bibitem[\citeproctext]{ref-zheng2023mtbench}
Zheng, L., Chiang, W.-L., Sheng, Y., Zhuang, S., Wu, Z., Zhuang, Y.,
Lin, Z., Li, Z., Li, D., Xing, E. P., Zhang, H., Gonzalez, J. E., \&
Stoica, I. (2023). Judging {LLM}-as-a-judge with {MT-Bench} and chatbot
arena. \emph{Advances in Neural Information Processing Systems (NeurIPS)
Datasets and Benchmarks Track}. \url{https://arxiv.org/abs/2306.05685}

\end{CSLReferences}

\end{document}
